\documentclass[a4paper,11pt,twoside,openright,numbers=noenddot]{scrreprt}

\makeatletter
%% ============================ SETTINGS ============================
%% Flip these three, then compile with pdflatex (twice) or lualatex.
\newif\ifhv@review  \hv@reviewfalse   % DRAFT watermark + visible \todo notes
\newif\ifhv@dutch   \hv@dutchfalse    % Dutch as main language
\newif\ifhv@secpage \hv@secpagetrue   % every \section starts on a new page
%% ==================================================================
% ---------------------------------------------------------------- engine/fonts
\usepackage{iftex}
\ifPDFTeX
  \usepackage[T1]{fontenc}
  \usepackage[utf8]{inputenc}
  \usepackage{textcomp}
  \usepackage{lmodern}
  \usepackage[scaled=0.94]{helvet}   % <- swap for the brand face
  \usepackage{courier}
\else
  \usepackage{fontspec}
  % \setsansfont{Your Brand Sans}[Scale=0.96]
  % \setmonofont{Your Brand Mono}[Scale=0.90]
\fi
\renewcommand{\familydefault}{\sfdefault}   % sans body text, typical for manuals

% ------------------------------------------------------------------- languages
\ifhv@dutch
  \usepackage[main=dutch,english]{babel}
\else
  \usepackage[main=english,dutch]{babel}
\fi
\usepackage{csquotes}
\usepackage{microtype}

% -------------------------------------------------------------------- geometry
\usepackage[a4paper,
  inner=28mm, outer=20mm, top=26mm, bottom=24mm,
  headsep=8mm, footskip=12mm]{geometry}

% ---------------------------------------------------------------------- colour
% Replace these three with the real HydroVolta brand values.
\usepackage[table]{xcolor}
\definecolor{hvprimary}{HTML}{003058}   % Hydrovolta navy
\definecolor{hvaccent} {HTML}{F0D828}   % Hydrovolta yellow (rules & fills only)
\definecolor{hvink}    {HTML}{17202A}
\definecolor{hvmuted}  {HTML}{6C7A89}
\definecolor{hvrule}   {HTML}{D3D9DF}
\definecolor{hvpaper}  {HTML}{F4F6F8}

% ISO 3864 / ANSI Z535 safety palette
% Red = personal injury (DANGER / WARNING / CAUTION)
% Yellow = property damage (NOTICE) · Green = information (NOTE)
\definecolor{hvdanger}   {HTML}{8A1414}
\definecolor{hvdangerbg}   {HTML}{F8E4E4}
\definecolor{hvwarning}  {HTML}{B01B1B}
\definecolor{hvwarningbg}  {HTML}{FBEDED}
\definecolor{hvcaution}  {HTML}{C24A4A}
\definecolor{hvcautionbg}  {HTML}{FDF2F2}
\definecolor{hvnotice}   {HTML}{F0D828}
\definecolor{hvnoticebg}   {HTML}{FEF9E0}
\definecolor{hvnoticeln} {HTML}{C9A21A}
\definecolor{hvnote}     {HTML}{3F6B4A}
\definecolor{hvnotebg}     {HTML}{EEF4EF}

% ------------------------------------------------------------------- packages
\usepackage{graphicx}
% ---- image folders -----------------------------------------------------------
% Project layout assumed:
%   HydroVolta-Main/
%     images/     <- shared image folder (sibling of this document's folder)
%     manual/     <- this LaTeX document
% Change BOTH lines below together if the sibling folder has another name.
\newcommand{\hvimgdir}{../images/}
\graphicspath{{../images/}{}}
\usepackage{tikz}
\usetikzlibrary{calc}
\usepackage[most]{tcolorbox}
\usepackage{booktabs,tabularx,longtable,xltabular,multirow,array,ragged2e}
\usepackage{enumitem}
\usepackage{siunitx}
\usepackage{listings}
\usepackage{qrcode}
\usepackage{lastpage}
\usepackage{caption}
\usepackage{pdfpages}
\usepackage{xspace}
\usepackage{scrlayer-scrpage}
\usepackage{etoolbox}
\usepackage{etoc}
\usepackage[hidelinks,bookmarksnumbered]{hyperref}
\usepackage[nameinlink]{cleveref}
\ifhv@review
  \usepackage{draftwatermark}
  \SetWatermarkText{DRAFT}\SetWatermarkScale{1.1}\SetWatermarkLightness{0.92}
  \usepackage[textwidth=18mm,textsize=tiny]{todonotes}
\else
  \usepackage[disable]{todonotes}
\fi

\sisetup{per-mode=symbol, range-units=single, detect-weight=true}

% ------------------------------------------------------------------- metadata
% Set every field from the main file with \hvsetup{key = value, ...}
\pgfkeys{/hv/.cd,
  title/.store in         = \hv@title,
  subtitle/.store in      = \hv@subtitle,
  product/.store in       = \hv@product,
  model/.store in         = \hv@model,
  serial/.store in        = \hv@serial,
  docnumber/.store in     = \hv@docnumber,
  revision/.store in      = \hv@revision,
  date/.store in          = \hv@date,
  language/.store in      = \hv@language,
  author/.store in        = \hv@author,
  approver/.store in      = \hv@approver,
  company/.store in       = \hv@company,
  address/.store in       = \hv@address,
  contact/.store in       = \hv@contact,
  website/.store in       = \hv@website,
  qrurl/.store in         = \hv@qrurl,
  logo/.store in          = \hv@logo,
  coverimage/.store in    = \hv@coverimage,
  classification/.store in= \hv@classification,
}
\newcommand{\hvsetup}[1]{\pgfqkeys{/hv}{#1}}
\hvsetup{title={}, subtitle={}, product={}, model={}, serial={},
  docnumber={}, revision={}, date={\today}, language={}, author={},
  approver={}, company={}, address={}, contact={}, website={}, qrurl={},
  logo={figures/logo.pdf}, coverimage={}, classification={}}

\AtBeginDocument{\hypersetup{pdftitle={\hv@title}, pdfauthor={\hv@company},
  pdfsubject={\hv@product\ \hv@model}, pdfkeywords={\hv@docnumber, \hv@revision}}}

% --------------------------------------------------------- headings & captions
\setkomafont{disposition}{\sffamily\bfseries\color{hvprimary}}
\setkomafont{chapter}{\sffamily\bfseries\huge\color{hvprimary}}
\setkomafont{section}{\sffamily\bfseries\Large\color{hvprimary}}
\setkomafont{subsection}{\sffamily\bfseries\large\color{hvink}}
\setkomafont{subsubsection}{\sffamily\bfseries\normalsize\color{hvink}}
\RedeclareSectionCommand[beforeskip=0pt,afterskip=20pt]{chapter}
\renewcommand*{\chapterformat}{}
\renewcommand{\chapterlinesformat}[3]{%
  \parbox{\linewidth}{%
    \textcolor{hvaccent}{\rule{28mm}{3pt}}\par\vspace{7pt}%
    \ifblank{#2}{}{{\normalfont\sffamily\bfseries\footnotesize\color{hvmuted}%
      \MakeUppercase{\chapapp}\;#2\par\vspace{5pt}}}%
    \raggedright #3\par}}
\setcounter{secnumdepth}{3}
% Optional one-line subtitle, placed directly after \chapter{...}
\newcommand{\chaptersubtitle}[1]{%
  \vspace{-12pt}\par
  {\sffamily\fontsize{13}{17}\selectfont\color{hvmuted}\raggedright #1\par}
  \vspace{10pt}}
\setcounter{tocdepth}{2}

\captionsetup{font={small},labelfont={bf,color=hvprimary},
  labelsep=period, justification=raggedright, singlelinecheck=false, skip=6pt}

% Every \section starts on a new page (disable with the nosectionpages option)
\ifhv@secpage\pretocmd{\section}{\clearpage}{}{%
  \PackageWarning{hydrovolta}{Could not patch \string\section}}\fi

% "In this chapter" list — place directly after \chapter / \chaptersubtitle
\newcommand{\chaptertoc}{%
  \begingroup
  \etocsetnexttocdepth{1}%
  \etocsettocstyle
    {{\sffamily\bfseries\footnotesize\color{hvmuted}IN THIS CHAPTER\par}%
     \vspace{5pt}{\color{hvink}\rule{\linewidth}{1.2pt}}\par\vspace{2pt}}{}%
  \etocsetstyle{section}{}{}
    {\par\vspace{4pt}\noindent
     \makebox[11mm][l]{\sffamily\bfseries\small\color{hvprimary}\etocnumber}%
     \sffamily\small\etocname\hfill{\color{hvmuted}\etocpage}%
     \par\vspace{4pt}{\color{hvrule}\rule{\linewidth}{0.4pt}}}{}%
  \localtableofcontents
  \endgroup}

\setlength{\parindent}{0pt}
\setlength{\parskip}{6pt plus 2pt}
\linespread{1.06}

% --------------------------------------------------------------- running heads
\pagestyle{scrheadings}
\clearpairofpagestyles
\automark[section]{chapter}
\setkomafont{pageheadfoot}{\normalfont\sffamily\footnotesize\color{hvmuted}}
\setkomafont{pagenumber}{\normalfont\sffamily\footnotesize\color{hvprimary}}
\ihead{\hv@product\ifx\hv@model\empty\else\ \textbar\ \hv@model\fi}
\ohead{\headmark}
\ofoot{\pagemark\,/\,\pageref{LastPage}}
\KOMAoptions{headsepline=0.4pt:\textwidth, footsepline=0.4pt:\textwidth}
\setkomafont{headsepline}{\color{hvrule}}
\setkomafont{footsepline}{\color{hvrule}}

% ------------------------------------------------------------- safety pictogram
% Simple ISO-style triangle; swap for the official symbol artwork when available.
\newcommand{\hvtriangle}[1][white]{%
  \raisebox{-0.6ex}{\begin{tikzpicture}[scale=0.10]
    \draw[line width=1.6pt, #1, line join=round] (0,0) -- (4,6.9) -- (8,0) -- cycle;
    \node[#1,font=\bfseries\tiny] at (4,2.0) {!};
  \end{tikzpicture}}}

% ---------------------------------------------------------------- admonitions
\tcbset{hvbox/.style={
  enhanced, breakable, sharp corners, boxrule=0.5pt,
  fonttitle=\sffamily\bfseries\small, coltitle=white,
  left=8pt, right=8pt, top=7pt, bottom=7pt,
  toptitle=3pt, bottomtitle=3pt,
  before skip=10pt, after skip=10pt}}

\newtcolorbox{danger} [1][]{hvbox, colframe=hvdanger,  colback=hvdangerbg,  colbacktitle=hvdanger,  title={\hvtriangle\ \MakeUppercase{danger}}, #1}
\newtcolorbox{warning}[1][]{hvbox, colframe=hvwarning, colback=hvwarningbg, colbacktitle=hvwarning, title={\hvtriangle\ \MakeUppercase{warning}}, #1}
\newtcolorbox{caution}[1][]{hvbox, colframe=hvcaution, colback=hvcautionbg, colbacktitle=hvcaution, title={\hvtriangle\ \MakeUppercase{caution}}, #1}
\newtcolorbox{notice} [1][]{hvbox, colframe=hvnoticeln, colback=hvnoticebg, colbacktitle=hvnotice, coltitle=hvink, title={\hvtriangle[hvink]\ \MakeUppercase{notice}}, #1}
\newtcolorbox{note}   [1][]{hvbox, colframe=hvnote,    colback=hvnotebg,    colbacktitle=hvnote,    title={\MakeUppercase{note}}, #1}

% -------------------------------------------------------------- lists & steps
\setlist[itemize]{leftmargin=1.2em, itemsep=2pt, topsep=4pt}
\setlist[enumerate]{leftmargin=1.6em, itemsep=2pt, topsep=4pt}
\newlist{procedure}{enumerate}{2}
\setlist[procedure,1]{label=\textbf{\arabic*.}, ref=\arabic*, leftmargin=2.0em,
  itemsep=8pt, topsep=8pt, parsep=3pt}
\setlist[procedure,2]{label=\alph*., leftmargin=1.6em, itemsep=3pt}
% Expected result of a step
\newcommand{\result}[1]{\par\smallskip{\color{hvmuted}$\rightarrow$\ \itshape #1}\par}

% Reference to an HMI element / physical control
\newcommand{\ui}[1]{\textbf{\color{hvprimary}#1}}
% Reference to a PLC tag / variable / signal name
\newcommand{\plctag}[1]{\texttt{\small #1}}
% Product names
\newcommand{\salinbloc}{SalinBloc\texttrademark\xspace}
\newcommand{\sonixed}{SonixED\texttrademark\xspace}
% Numbered balloon for figure callouts
\newcommand{\callout}[1]{\tikz[baseline=(n.base)]{\node[circle,fill=hvprimary,
  text=white,inner sep=1.4pt,font=\sffamily\bfseries\tiny,minimum size=11pt](n){#1};}}

% ---------------------------------------------------------------------- tables
% \hvtablebegin{colspec}{caption}{label} ... rows ... \bottomrule\end{xltabular}\par
% NOTE: xltabular/tabularx locate their own closing \end{xltabular} by
% scanning the raw input for the *first* literal "\end{...}" token. That
% scan cannot be hidden inside another environment or macro — doing so
% (e.g. \newenvironment{hvtable}{...}{\end{xltabular}} or an \end-hiding
% \newcommand) makes the scanner swallow the wrapper's own closing tag
% instead, so it never finds \end{xltabular} and runs off the end of the
% file ("File ended while scanning use of \TX@get@body."). So every use
% must end with the literal "\bottomrule\end{xltabular}\par" typed directly
% — only the opening half is safe to wrap in a macro.
\newcommand{\hvtablebegin}[3]{\par\vspace{10pt}\captionof{table}{#2}\label{#3}\par\vspace{3pt}%
   \small\setlength{\tabcolsep}{6pt}\renewcommand{\arraystretch}{1.25}%
   \begin{xltabular}{\linewidth}{#1}\toprule}
\newcommand{\hvhead}[1]{\rowcolor{white}#1}
\newcolumntype{L}{>{\RaggedRight\arraybackslash}X}
\newcolumntype{C}[1]{>{\Centering\arraybackslash}p{#1}}
\newcolumntype{R}[1]{>{\RaggedLeft\arraybackslash}p{#1}}

% Abbreviations / glossary block:
% \abbreviationsbegin ... \abbr{term}{def} ... \bottomrule\end{xltabular}\par
\newcommand{\abbreviationsbegin}{\par\small\renewcommand{\arraystretch}{1.3}%
   \begin{xltabular}{\linewidth}{C{28mm} L}\toprule
   \textbf{Term} & \textbf{Definition}\\\midrule}
\newcommand{\abbr}[2]{#1 & #2\\}

% ---------------------------------------------------------------------- figures
% \hvfig[width]{file}{caption}{label} - draws a placeholder if the file is absent
\newcommand{\hvplaceholder}[2][0.75\linewidth]{%
  \begin{tikzpicture}
    \node[draw=hvrule, dashed, fill=hvpaper, minimum width=#1,
          minimum height=48mm, align=center, inner sep=8pt,
          font=\ttfamily\scriptsize, text=hvmuted] {#2};
  \end{tikzpicture}}
% Resolve an image name against ./ , ./figures/ and \hvimgdir ; result in \hv@found
\newcommand{\hv@resolve}[1]{%
  \def\hv@found{}%
  \IfFileExists{#1}{\def\hv@found{#1}}{%
    \IfFileExists{figures/#1}{\def\hv@found{figures/#1}}{%
      \IfFileExists{\hvimgdir#1}{\edef\hv@found{\hvimgdir#1}}{}}}}
\newcommand{\hvfig}[4][0.75\linewidth]{%
  \begin{figure}[htbp]\centering
    \hv@resolve{#2}%
    \ifx\hv@found\@empty\hvplaceholder[#1]{#2}\else
      \includegraphics[width=#1]{\hv@found}\fi
    \caption{#3}\label{#4}
  \end{figure}}

% ------------------------------------------------------------------- QR module
\newcommand{\hvqr}[2][22mm]{\qrcode[height=#1,level=M]{#2}}
\newcommand{\hvqrbox}[3][22mm]{%
  \begin{tcolorbox}[hvbox, colframe=hvrule, colback=hvpaper]
    \begin{minipage}[c]{#1}\hvqr[#1]{#2}\end{minipage}\hfill
    \begin{minipage}[c]{\dimexpr\linewidth-#1-10pt\relax}\small #3\end{minipage}
  \end{tcolorbox}}

% ---------------------------------------------- IEC 61131-3 Structured Text
\lstdefinelanguage{ST}{
  morekeywords={PROGRAM,END_PROGRAM,FUNCTION,END_FUNCTION,FUNCTION_BLOCK,
    END_FUNCTION_BLOCK,VAR,VAR_INPUT,VAR_OUTPUT,VAR_IN_OUT,VAR_GLOBAL,VAR_TEMP,
    END_VAR,CONSTANT,RETAIN,PERSISTENT,AT,TYPE,END_TYPE,STRUCT,END_STRUCT,ARRAY,
    OF,IF,THEN,ELSIF,ELSE,END_IF,CASE,END_CASE,FOR,TO,BY,DO,END_FOR,WHILE,
    END_WHILE,REPEAT,UNTIL,END_REPEAT,EXIT,CONTINUE,RETURN,TRUE,FALSE,AND,OR,
    XOR,NOT,MOD,ACTION,END_ACTION,STEP,INITIAL_STEP,END_STEP,TRANSITION,
    END_TRANSITION,FROM,METHOD,END_METHOD,INTERFACE,END_INTERFACE,EXTENDS,
    IMPLEMENTS,CONFIGURATION,END_CONFIGURATION,RESOURCE,END_RESOURCE,TASK,ON,
    WITH,PRIORITY,INTERVAL,SINGLE},
  morekeywords=[2]{BOOL,BYTE,WORD,DWORD,LWORD,SINT,INT,DINT,LINT,USINT,UINT,
    UDINT,ULINT,REAL,LREAL,TIME,LTIME,DATE,TIME_OF_DAY,TOD,DATE_AND_TIME,DT,
    STRING,WSTRING,POINTER,REFERENCE,ANY,ANY_INT,ANY_REAL},
  sensitive=false,
  morecomment=[l]{//},
  morecomment=[s]{(*}{*)},
  morecomment=[s]{/*}{*/},
  morestring=[b]',
  morestring=[b]"}

\lstdefinestyle{hvcode}{
  language=ST,
  basicstyle=\ttfamily\scriptsize\color{hvink},
  keywordstyle=\color{hvprimary}\bfseries,
  keywordstyle=[2]\color{hvaccent},
  commentstyle=\color{hvmuted}\itshape,
  stringstyle=\color{hvnote},
  numbers=left, numberstyle=\ttfamily\tiny\color{hvmuted}, numbersep=8pt,
  backgroundcolor=\color{hvpaper},
  frame=leftline, framerule=2pt, rulecolor=\color{hvrule},
  framesep=8pt, xleftmargin=16pt,
  showstringspaces=false, breaklines=true, breakatwhitespace=true,
  tabsize=2, captionpos=b, aboveskip=10pt, belowskip=10pt}
\lstset{style=hvcode}
\newcommand{\hvcodefile}[3]{\lstinputlisting[caption={#2},label={#3}]{#1}}
\renewcommand{\lstlistingname}{Listing}

% --------------------------------------------------------------- cover page
\newcommand{\hvtitlepage}{%
  \begin{titlepage}
  \thispagestyle{empty}
  \begin{tikzpicture}[remember picture, overlay]
    \fill[hvprimary] (current page.north west) rectangle
      ($(current page.north east)+(0,-52mm)$);
    \fill[hvaccent] ($(current page.north west)+(0,-52mm)$) rectangle
      ($(current page.north east)+(0,-55mm)$);
    \node[anchor=north west, inner sep=0pt]
      at ($(current page.north west)+(20mm,-18mm)$)
      {\hv@resolve{\hv@logo}%
       \ifx\hv@found\@empty\color{white}\sffamily\Large\bfseries\hv@company
       \else\includegraphics[height=12mm]{\hv@found}\fi};
    \ifx\hv@classification\empty\else
      \node[anchor=north east, font=\sffamily\scriptsize, text=white!80!hvprimary]
        at ($(current page.north east)+(-20mm,-18mm)$) {\MakeUppercase{\hv@classification}};
    \fi
  \end{tikzpicture}

  \vspace*{70mm}
  {\sffamily\bfseries\fontsize{36}{42}\selectfont\color{hvprimary}\hv@title\par}
  \ifx\hv@subtitle\empty\else
    \vspace{6mm}{\sffamily\fontsize{15}{20}\selectfont\color{hvmuted}\hv@subtitle\par}
  \fi
  \vspace{9mm}
  {\color{hvaccent}\rule{28mm}{3pt}}\par
  \vspace{5mm}
  {\sffamily\large\color{hvink}\hv@product\ifx\hv@model\empty\else\ \textbar\ \hv@model\fi\par}

  \vfill
  \ifx\hv@coverimage\empty\else
    \begin{center}
      \hv@resolve{\hv@coverimage}%
      \ifx\hv@found\@empty\hvplaceholder[0.9\linewidth]{\hv@coverimage}%
      \else\includegraphics[width=0.9\linewidth]{\hv@found}\fi
    \end{center}
    \vfill
  \fi

  \end{titlepage}}

% ------------------------------------------------------- document control page
\newcommand{\hvdoccontrol}{%
  \thispagestyle{scrheadings}
  \section*{Document control}
  \small
  \begin{tabularx}{\linewidth}{@{}p{34mm} L@{}}
    \toprule
    Title            & \hv@title\\
    Product          & \hv@product\ifx\hv@model\empty\else\ \textbar\ \hv@model\fi\\
    Date of issue    & \hv@date\\
    Classification   & \hv@classification\\
    \bottomrule
  \end{tabularx}
  \par\vspace{10mm}
  \subsection*{Manufacturer}
  \small
  \hv@company\par\hv@address\par\hv@contact\par
  \ifx\hv@website\empty\else\href{\hv@website}{\hv@website}\par\fi
  \clearpage}

% ------------------------------------------------------------------ back cover
\newcommand{\hvbackcover}{%
  \clearpage\thispagestyle{empty}
  \begin{tikzpicture}[remember picture, overlay]
    \fill[hvprimary] ($(current page.south west)$) rectangle
      ($(current page.south east)+(0,54mm)$);
    \node[anchor=south west, align=left, text=white, font=\sffamily\small]
      at ($(current page.south west)+(20mm,16mm)$)
      {\textbf{\hv@company}\\[2pt]\hv@address\\[2pt]\hv@contact\\[2pt]\hv@website};
    \node[anchor=south east, align=right, text=white!80!hvprimary,
          font=\sffamily\scriptsize]
      at ($(current page.south east)+(-20mm,16mm)$)
      {\hv@docnumber\ \textbar\ Rev.\ \hv@revision\ \textbar\ \hv@date};
  \end{tikzpicture}
  \null}
\makeatother

%% ========================= DOCUMENT METADATA =========================
\hvsetup{
  title          = {},
  subtitle       = {},
  product        = {SalinBloc\texttrademark},
  model          = {},
  docnumber      = {},
  revision       = {},
  date           = {\today},
  language       = {English},
  author         = {},
  approver       = {},
  company        = {Hydrovolta NV},
  address        = {Romeinse Straat 18, 3001 Leuven, Belgium},
  contact        = {info@hydrovolta.com \textbar\ +32 28 800 336},
  website        = {https://hydrovolta.com},
  qrurl          = {https://hydrovolta.com/salinbloc-manual}, % URL behind the QR code
  logo           = {figures/logo-white.png},
  coverimage     = {figures/salinbloc-skid.png},
  classification = {},
}

%% Image folders — \hvfig{pump.png}{...}{fig:pump} needs no path.
%% Change both lines together if the shared folder has another name:
% \renewcommand{\hvimgdir}{../graphics/}
% \graphicspath{{../graphics/}{figures/}{}}

%%%%%%%%%%%%%%%%%%%%%%%%%%%%%%%%%%%%%%%%%%%%%%%%%%%%%%%%%%%%%%%%%%%%%%%%%%%%%%%
\begin{document}

\hvtitlepage

\pagenumbering{roman}
\hvdoccontrol
\tableofcontents
\clearpage
% \listoffigures
% \listoftables

\cleardoublepage\pagenumbering{arabic}\setcounter{page}{1}

%%%%%%%%%%%%%%%%%%%%%%%%%%%%%%%%% CONTENT %%%%%%%%%%%%%%%%%%%%%%%%%%%%%%%%%%%%%

%% ---- Part I: General chapters ------------------------------------------

\chapter{About this manual}
\label{ch:about}
\chaptertoc

\section{Purpose and scope}
This manual describes the \salinbloc container: a modular \sonixed
electrodialysis unit for water desalination and brine valorization,
delivered as one or more standard shipping containers fitted with
pretreatment, stack, dosing and control equipment. Part~I of this manual
covers the unit as a whole (safety, installation, commissioning, day-to-day
operation). Part~II gives one self-contained chapter per compartment inside
the container --- pumps, stacks, valves, dosing, sensors, control cabinet,
tanks, and telemetry --- each with its own specifications, operating notes,
maintenance and troubleshooting, so it can be consulted (or printed/exported)
on its own. Part~III closes with shutdown/storage and reference appendices.

\sonixed is HydroVolta NV's electrodialysis-reversal (EDR) technology,
combined with an integrated ultrasound antifouling stage and a custom
ion-selective membrane coating. Hydrovolta NV (Leuven, Belgium, founded
2016) developed \sonixed over eight years of lab validation with research
partners including VITO, and holds a granted European patent (2021) on the
combined EDR~+ ultrasound process. \sonixed/\salinbloc units are in active
use or pilot deployment at sites in Belgium, Morocco, Australia (Perth),
Brazil and Saudi Arabia; see \cref{sec:fieldperformance} for field
performance data published by HydroVolta.

\begin{notice}
Several sections in this manual are marked \textbf{TODO}. These mark
information that was not found in HydroVolta's engineering documentation or
public materials at the time of writing (July--August~2026) and must be
supplied by HydroVolta engineering, or confirmed by inspection/measurement
on an actual unit, before this manual is issued to a customer. Every other
statement in this manual is grounded in a specific HydroVolta engineering
document, datasheet, or the hydrovolta.com website --- but the underlying
documents mostly describe pilot/demonstration containers (Farys, Ganzepoot,
BASF, Morocco, and others), not a finalized, branded \enquote{\salinbloc}
bill of materials. Treat pilot-specific facts as a well-informed engineering
basis for the standard product, not as a guaranteed as-shipped
specification, until confirmed against a specific unit's documentation.
Notably, the \textbf{PLC control program itself is out of scope for this
edition} of the manual --- it is still under development by HydroVolta's
automation team, and the relevant sections (HMI screen navigation, exact
control logic, full setpoint list) will be completed once that work is
finished.
\end{notice}

\section{Who this manual is for}
This manual is written for the people who take delivery of, install,
commission, operate and maintain a \salinbloc unit: site engineers,
installation contractors, plant operators, and HydroVolta service
technicians. Readers are assumed to have basic familiarity with industrial
piping, low-voltage electrical systems and PLC-controlled equipment; no
prior knowledge of electrodialysis is assumed (see the Glossary,
\cref{ch:appglossary}, for unfamiliar terms).

\section{Symbols and signal words}
Hazard and information boxes in this manual follow the signal-word
convention below, aligned with ISO~3864 / ANSI~Z535 safety-color practice.

\begin{danger}
Indicates a hazard that, if not avoided, \textbf{will} result in death or
serious injury.
\end{danger}
\begin{warning}
Indicates a hazard that, if not avoided, \textbf{could} result in death or
serious injury.
\end{warning}
\begin{caution}
Indicates a hazard that, if not avoided, could result in minor or moderate
injury.
\end{caution}
\begin{notice}
Indicates a risk of property damage or equipment malfunction, with no
personal-injury hazard.
\end{notice}
\begin{note}
Highlights additional information that helps the reader but is not itself
a hazard warning.
\end{note}

\section{Online documentation and QR code}
A QR code on the container's control cabinet links to the latest online
version of this manual and to the unit's factory quality report.

\makeatletter
\hvqrbox{\hv@qrurl}{Scan to open the online manual and this unit's
  quality report at \hv@qrurl.}
\makeatother

\section{Contact and support}
Hydrovolta NV, Romeinse Straat~18, 3001 Leuven, Belgium.
E-mail:~\href{mailto:info@hydrovolta.com}{info@hydrovolta.com}. Phone:
+32~28~800~336. Web:~\href{https://hydrovolta.com}{hydrovolta.com}. TODO:
no dedicated \emph{service} phone line, out-of-hours escalation path, or
service-hours commitment (as opposed to general company contact) was
found --- see also \cref{sec:whentocallservice} for when to contact
service.


\chapter{Safety}
\label{ch:safety}
\chaptertoc

\begin{warning}
Read this entire chapter before installing, operating or servicing the
\salinbloc unit. No unit-level hazard analysis, PPE schedule or
lockout/tagout procedure specific to \salinbloc/\sonixed exists yet in
HydroVolta's engineering library --- only generic contractor-safety
material (tool safety sheets, an electrical risk assessment template, a
VCA contractor-competency handbook) and the Safety Data Sheets for the
dosing/CIP chemicals. The hazards, PPE and isolation guidance below are
derived from those chemical SDSs and from the electrical/pressure/DC
figures documented elsewhere in this manual (\cref{ch:cabinet,ch:stacks});
they are a defensible starting point, \textbf{not} a substitute for a
formal hazard analysis, which HydroVolta engineering/EHS must still
produce and sign off before this manual is issued to a customer.
\end{warning}

\section{Intended use}
\label{sec:intendeduse}
The \salinbloc unit is intended for the electrodialysis desalination or ion
reduction of brackish groundwater, industrial process water, or brine, in
the total dissolved solids range HydroVolta publishes as the technology's
operating envelope --- approximately 1{,}500--20{,}000\,ppm TDS (see
\cref{sec:fieldperformance}) --- and within the feedwater quality limits
stated for the pretreatment and stack stages
(\cref{ch:pretreatment,ch:stacks}). It is not intended for seawater
desalination (HydroVolta states seawater, $>$35{,}000\,ppm, as outside
\sonixed's primary use case) and is \textbf{not} currently certified as a
drinking-water treatment system: the \texttt{03\_UK DWI reg~31 approval}
folder in HydroVolta's datasheet library contains only generic
third-party guidance on how Reg~31 approval works (not an actual approval
or application), and the \texttt{04\_NSF} folder contains only a general
NSF/ANSI~60/61/372 adoption survey (not an NSF listing) --- confirmed by
direct inspection of both folders. Do not make a potable-water regulatory
claim to a customer until formal approval is obtained --- see
\cref{ch:appcert}. Foreseeable misuse includes: feeding
water outside the documented turbidity/iron/chlorine/pH envelope
(\cref{ch:stacks}) without adequate pretreatment; bypassing or defeating
interlocks, alarms or the emergency stop; operating with the cabinet doors
open; and mixing CIP chemicals outside the sequence in
\cref{ch:dosing} (in particular, never combine NaOH with feed/process
water --- see the warning in \cref{ch:dosing}).

\section{Qualified personnel}
Work on the control cabinet, VFDs and any live electrical circuit,
including the \sonixed stack's DC bus (up to 600\,V\,DC,
\cref{ch:stacks}), must only be performed by an electrically qualified
person as defined in IEC~60204-1. Chemical dosing tank refilling and CIP
chemical handling must only be performed by personnel trained on the
relevant Safety Data Sheet for that chemical (\cref{sec:mainhazards}).
HydroVolta's own engineering library uses the Belgian/Dutch \textbf{VCA}
(Veiligheid, Gezondheid en Milieu Checklist Aannemers) contractor-safety
certification as its baseline competency reference for site work; treat
VCA-certified (or locally equivalent) status as the expected minimum for
any contractor performing installation or maintenance work on the unit.
General operators (start/stop, HMI monitoring, routine visual checks) do
not need electrical or chemical-handling qualifications but must be
trained on this manual and on the local isolation procedure
(\cref{sec:isolation}) before working unsupervised.

\section{Personal protective equipment}
\label{sec:ppe}
\begin{itemize}
\item \textbf{Chemical dosing / CIP work} (antiscalant, HCl, NaOH, NaOCl,
  H\textsubscript{2}SO\textsubscript{4} --- all corrosive per their SDS,
  \cref{sec:mainhazards}): chemical-resistant gloves, safety
  glasses/goggles with a face shield for tank filling or CIP setup, and a
  chemical-resistant apron. Have an eyewash/emergency shower available
  wherever CIP or dosing-tank work is performed.
\item \textbf{Control cabinet / electrical work}: PPE appropriate to the
  task and the incoming supply voltage (3$\times$400\,V\,AC,
  \cref{ch:cabinet}) per the site's local electrical safety regulation and
  the qualified person's own risk assessment; always verify isolation
  (\cref{sec:isolation}) before opening the cabinet.
\item \textbf{General site / mechanical work} (pumps, valves, piping):
  safety shoes and, where local site rules require it, hi-vis clothing and
  hearing protection near running pumps.
\end{itemize}
TODO: this list is derived from the chemical SDSs and general industrial
practice, not from a HydroVolta-authored PPE matrix --- none was found.

\section{Main hazards}
\label{sec:mainhazards}
\begin{itemize}
\item \textbf{Electrical:} the incoming 3$\times$400\,V\,AC cabinet supply
  (\cref{ch:cabinet}), and --- specific to this technology --- the
  \sonixed stack's own DC bus, which is energized to \textbf{60--600\,V\,DC
  at 3--20\,A} during normal operation (\cref{ch:stacks}). This DC bus is
  a hazard independent of the AC supply and must be isolated and verified
  dead separately before any stack or manifold work.
\item \textbf{Chemical:} the antiscalant, HCl, NaOH, NaOCl and
  H\textsubscript{2}SO\textsubscript{4} used for dosing and CIP
  (\cref{ch:dosing}) are corrosive per their Safety Data Sheets (Carl
  Roth SDS on file in HydroVolta's engineering library). \textbf{Never mix
  NaOH with feed or process water} --- HydroVolta's own CIP procedure
  documents this as a scaling hazard, and mixing acid and hypochlorite
  streams carries a toxic-gas risk common to all such systems; keep
  chemical dosing lines and procedures strictly separated as described in
  \cref{ch:dosing}.
\item \textbf{Pressurized piping:} the hydraulic circuit operates up to
  2.5\,bar (stack maximum 3\,bar, \cref{ch:stacks}) and the pretreatment
  UF stage's CEB/backwash cycle briefly reaches similar pressures
  (\cref{ch:pretreatment}); depressurize before opening any fitting.
\item \textbf{Moving parts:} process pumps, dosing pumps, and automated
  valve actuators (\cref{ch:pumps,ch:valves}) start automatically under
  PLC control without warning --- treat every pump and valve as capable of
  starting at any time unless the unit is isolated.
\end{itemize}

\hvfig{fig-safety-hazard-icons.png}{Suggested figure: a labeled diagram of
  the container's exterior and interior showing each hazard zone above
  (electrical cabinet / DC stack bus, chemical dosing area, pressurized
  piping, pump/valve zones) with its corresponding ISO~7010 warning
  pictogram.}
  {fig:safety-hazards}

\section{First aid and chemical spill response}
\label{sec:firstaid}
The following is transcribed directly from the manufacturer Safety Data
Sheets on file for the dosing/CIP chemicals (Carl Roth SDS, Section~4
\enquote{First aid measures} and Section~6 \enquote{Accidental release
measures}; text is consistent across NaOH, HCl, NaOCl and
H\textsubscript{2}SO\textsubscript{4}). Post this section, or an
equivalent quick-reference card, wherever chemicals are handled.

\begin{danger}
\textbf{Eye contact:} flush immediately with plenty of flowing water for
\textbf{10 to 15 minutes}, holding eyelids apart, and consult an
ophthalmologist. Protect the uninjured eye. \textbf{Skin contact:} take
off all contaminated clothing immediately; wash the affected skin
immediately with plenty of water (for H\textsubscript{2}SO\textsubscript{4}:
first mechanically remove the chemical, e.g.\ dab away with wadding, then
wash with a mild cleansing agent and water); seek immediate medical
treatment --- corrosive injuries left untreated are hard to cure.
\textbf{Inhalation:} move to fresh air; seek medical advice if symptoms
persist or in any case of doubt. \textbf{Ingestion:} rinse the mouth
immediately and drink plenty of water; call a physician immediately ---
risk of gastric/esophageal perforation from the corrosive effect.
\end{danger}

\textbf{Spill response} (non-emergency personnel only; do not attempt
without the PPE in \cref{sec:ppe}): put on PPE, avoid contact with skin,
eyes and clothing, and do not breathe vapour/spray. Cover nearby drains to
protect surface/ground water --- for HCl and H\textsubscript{2}SO\textsubscript{4}
specifically, the product must be neutralized before any discharge to a
drain or sewage plant. Absorb the spill with a liquid-binding material
(sand, diatomaceous earth, or an acid/universal binding agent), then place
the absorbed material in an appropriate container for disposal. Ventilate
the affected area, especially after an HCl release.

\section{Safety devices and emergency stop}
HydroVolta's HMI design specification (\cref{ch:operation}) confirms a
\textbf{hardware Emergency Stop button that bypasses the PLC entirely} ---
see \cref{sec:estop} for the full recovery procedure. TODO: its physical
location on the control panel/cabinet, and whether more than one E-stop
station exists (e.g.\ one per container, or one per access point), was
not confirmed in the documents reviewed --- verify against the as-built
electrical schematic and IEC~60204-1 on the actual unit before
publication. No document describing pressure-relief devices or leak
detection was found.

\section{Isolating the unit before work}
\label{sec:isolation}
No HydroVolta-authored lockout/tagout procedure was found; until one
exists, apply this general sequence before any maintenance, sensor, or
piping work:
\begin{procedure}
\item Stop the process via the normal \plctag{Stop} sequence
  (\cref{ch:operation}).
\item Isolate and lock out the incoming electrical supply at the main
  breaker; verify absence of voltage at the cabinet terminals before
  touching any wiring.
\item Separately verify the \sonixed stack's DC bus is de-energized ---
  it operates independently of the AC supply isolation point and can
  present up to 600\,V\,DC (\cref{ch:stacks}); never assume it is safe
  just because the AC breaker is open.
\item Isolate and depressurize the hydraulic circuit (close feed/product/
  drain isolation points, \cref{sec:watercx}) before opening any fitting.
\item Where chemical dosing lines will be disturbed, close the relevant
  tank isolation valve and follow the PPE guidance in \cref{sec:ppe}.
\end{procedure}
TODO: this sequence is assembled from the general electrical/hydraulic/DC
facts documented elsewhere in this manual, not from a validated
HydroVolta LOTO procedure --- HydroVolta engineering/EHS must review and
formalize it, including lock/tag hardware and a permit-to-work step if
required by the customer's site rules.

\section{Safety labels on the container}
A folder named \texttt{17\_Photos/Container stickers} exists in
HydroVolta's engineering library, but on inspection it contains only
\textbf{brand artwork} (the Hydro-Blue and Volta-Yellow logo marks, the
HydroVolta wordmark, a \enquote{water desalination} decal, and a door
logo) --- \textbf{no hazard/warning pictogram stickers were found there or
anywhere else in the library}. This is a real gap: before a unit ships, it
needs a proper ISO~7010 hazard-label set (electrical, DC stack, chemical,
pressure --- matching \cref{sec:mainhazards}) actually affixed and
photographed, and this section rewritten from those photographs.

\hvfig{fig-safety-labels.png}{Suggested figure: photograph(s) of the
  container exterior and control cabinet with each safety/identification
  label numbered, plus a key table listing what each label means and where
  it is mounted --- to be produced once a real hazard-label set has been
  designed and affixed to a unit (none currently exists in HydroVolta's
  documentation).}
  {fig:safety-labels}


\chapter{System overview}
\label{ch:overview}
\chaptertoc

\section{What the unit does}
The \salinbloc unit is built around HydroVolta's \sonixed technology,
which combines three mechanisms in one automated platform:
\begin{enumerate}
\item \textbf{Electrodialysis reversal (EDR):} an applied DC electric
  field draws dissolved ions --- not water --- through ion-exchange
  membranes from the diluate stream into the concentrate stream. Because
  water itself does not have to be forced through a membrane (unlike
  reverse osmosis), the process runs at much lower hydraulic pressure.
  Periodically reversing polarity cleans the membranes in place without
  chemicals and gives more stable performance as feed salinity varies.
\item \textbf{Integrated ultrasound antifouling:} ultrasound transducers
  built into the stack create acoustic cavitation at the membrane
  surface, disrupting scale nucleation and helping prevent silica
  fouling --- this is the element of \sonixed that HydroVolta holds a
  granted European patent on (2021).
\item \textbf{Custom membrane coating:} a proprietary coating gives the
  membranes monovalent-ion selectivity, so the process can target
  nitrate, sodium, chloride and fluoride removal while retaining calcium
  and magnesium in the product water.
\end{enumerate}
The result is a saline feed stream split into a low-salinity product
stream (\emph{diluate}) and a concentrated reject stream
(\emph{concentrate}/brine), with chemical-free ultrasound antifouling
keeping the membrane stack clean and extending electrode/membrane life
(\cref{ch:stacks}). Typical applications: groundwater desalination and
selective ion (e.g.\ nitrate) removal, industrial process-water reuse,
brine valorization/concentration, and Zero Liquid Discharge (ZLD) support.

\hvfig{fig-container-exterior.png}{Suggested figure: exterior photograph of
  a delivered \salinbloc container (or container pair) on site, showing
  the overall footprint and connection points.}{fig:container-exterior}

\section{Field performance reference}
\label{sec:fieldperformance}
HydroVolta publishes the following performance comparison from a 90-day
industrial pilot (February--May~2023, cooling-tower blowdown at a chemical
industry facility in Germany, feed TDS $\approx$4{,}500\,ppm, silica
85\,mg/l, high hardness), run head-to-head against reverse osmosis (RO)
and a standard EDR system on the same feedwater (source:
hydrovolta.com/technology.html and /projects.html). These figures describe
that specific pilot and HydroVolta's general \sonixed technology, not a
performance guarantee for any individual \salinbloc unit; use them as
context for what the technology is designed to achieve, and always verify
against the site-specific commissioning data recorded in
\cref{ch:appparams}.

\hvtablebegin{L R{18mm} R{18mm} R{22mm}}{Germany 2023 industrial pilot --- \sonixed vs.\ RO vs.\ standard EDR (source: hydrovolta.com)}{tab:germanypilot}
\hvhead{\textbf{Parameter} & \textbf{\sonixed} & \textbf{RO} & \textbf{Standard EDR}}\\\midrule
Water recovery & 90--98\% & 50--75\% & 70--85\%\\
Energy consumption & 0.19--0.7\,kWh/m\textsuperscript{3} & 2.0--4.0\,kWh/m\textsuperscript{3} & 1.0--2.0\,kWh/m\textsuperscript{3}\\
Operating pressure & 0.5--2.5\,bar & 25--35\,bar & 1--2\,bar\\
Chemical dosing & Minimal, no antiscalant & Required & Antiscalant needed\\
CIP frequency & 0.6$\times$/month & 1$\times$/month & 1$\times$/month\\
Total treatment cost & \$0.30--0.40/m\textsuperscript{3} & \$0.60--1.00/m\textsuperscript{3} & \$0.40--0.80/m\textsuperscript{3}\\
\bottomrule\end{xltabular}\par

HydroVolta states \sonixed's general operating envelope as feed TDS
1{,}500--20{,}000\,ppm (\enquote{strong fit} also for high silica
$>$30\,mg/l and high hardness $>$300\,mg/l Ca+Mg feeds); seawater
($>$35{,}000\,ppm) is stated as not the primary use case. Active and
piloted \sonixed/\salinbloc deployments referenced by HydroVolta include a
flagship groundwater desalination project in Morocco (seawater-intruded
coastal aquifer, $\approx$7{,}000\,ppm feed), a multi-site utility pilot
with Water Corporation in Perth, Australia (selective nitrate and
salinity removal), the Germany industrial pilot above, and earlier-stage
projects in Brazil and Saudi Arabia.

\section{Container layout}
The system is delivered as a modular set of skids assembled inside one or
more standard shipping containers (20\,ft and 40\,ft, High-Cube and
Standard variants). Larger installations use \textbf{two} containers: a
\emph{pretreatment container} (filtration stages, \cref{ch:pretreatment})
and a \emph{treatment container} housing the \sonixed skid(s)
(\cref{ch:pumps,ch:stacks,ch:valves,ch:cabinet}), connected hydraulically.
A documented 20\,ft general-arrangement layout (drawing 2308CA1) shows
internal dimensions of 6058\,$\times$\,2438\,$\times$\,2896\,mm, housing a
bank of \sonixed stack racks, an electrical/control cabinet, a
dosing/chemical section, three overhead storage tanks, and interconnecting
piping; a companion top-view drawing (2308CA2) lays the skids out along the
container floor: chemical skid, pump skid, tank skid, control~1 skid,
filter skid, \sonixed skid, control~2 skid.

\hvfig{fig-container-topview.png}{Suggested figure: the skid top-view
  layout (based on drawing 2308CA2), redrawn/simplified for this manual,
  with each skid numbered and cross-referenced to its Part~II chapter.}
  {fig:container-topview}

\section{Process flow}
\begin{procedure}
\item Feed water enters the pretreatment container (2\textonehalf" connection
  from source).
\item Pretreatment (disc filtration and ultrafiltration, \cref{ch:pretreatment})
  conditions the water; filtered water leaves via a 2" connection.
\item Filtered water enters the treatment container (2" connection) and is
  split and pumped through parallel \sonixed stacks
  (\cref{ch:pumps,ch:stacks}), with three internal streams per stack:
  \textbf{diluate} (becomes the product/permeate), \textbf{concentrate}
  (becomes the brine/reject), and \textbf{electrolyte} (a separate loop
  feeding the electrode compartments, \cref{ch:tanks}).
\item Product (permeate) water exits the treatment container via a 2"
  outlet connection.
\item Reject/concentrate and pretreatment backwash/drainage are combined
  via a tee piece and routed to a drain tank via 3" drainage connections
  from both containers.
\end{procedure}

\begin{note}
TODO: a fully tagged P\&ID with valve/instrument numbers exists for the
BASF pilot (drawing SWHC-001/00, rev.\,0.98) --- see \cref{ch:appdrawings}
--- but has not yet been confirmed as representative of the standard
\salinbloc product configuration.
\end{note}

\section{Main components}
This section is a one-page index of every compartment covered in its own
chapter later in this manual (Part~II) --- use it to jump straight to the
part you need.

\begin{itemize}
\item \textbf{Pretreatment and filtration skid} --- \cref{ch:pretreatment}
\item \textbf{Feed and process pumps} --- \cref{ch:pumps}
\item \textbf{\sonixed stack modules} --- \cref{ch:stacks}
\item \textbf{Valves and manifold} --- \cref{ch:valves}
\item \textbf{Dosing system} --- \cref{ch:dosing}
\item \textbf{Instrumentation and sensors} --- \cref{ch:sensors}
\item \textbf{Control cabinet, PLC and HMI} --- \cref{ch:cabinet}
\item \textbf{Tanks} --- \cref{ch:tanks}
\item \textbf{Telemetry and remote connectivity} --- \cref{ch:telemetry}
\end{itemize}

\section{Control cabinet and HMI}
The control system consists of a PLC with a touchscreen HMI, manual/override
controls, and RS485 communication, housed in an IP55-rated cabinet. The
skid ordering code allows two control-box configurations --- \textbf{CB0}
(PLC only) and \textbf{CB1} (PLC~+ touchscreen HMI) --- and three
cloud-connectivity tiers: \textbf{CL0} (local only), \textbf{CL1} (cloud
monitoring), \textbf{CL2} (full cloud monitoring and control). See
\cref{ch:cabinet} for details.

\section{Nameplate and identification}
TODO: no nameplate drawing, CE-marking template with filled-in fields, or
Declaration of Conformity content was found (the \texttt{02\_DoC's} and
\texttt{05\_SonixED stack} CE-marking folders in HydroVolta's engineering
library are confirmed empty --- see \cref{ch:appcert}). The confirmed
identification data is the datasheet labeling convention \enquote{Model:
SONIXED-SKID, Version: 2025-A7} and the ordering code format
\texttt{SONIXED-SKID-A-S-M-PR-V-CB-CL} (application type A0/A1/A2, stack
count 1/2/3, measuring points 6/9/12, pressure sensors 6/9/12, valve
count 2/4/6, control-box option CB0/CB1, cloud option CL0/CL1/CL2). Each
skid physically measures approximately 3040~$\times$~1340~$\times$~2020\,mm
(l~$\times$~b~$\times$~h), on an aluminum-profile frame with a
PVC/thermoplastic-spacer housing, mounted on lockable wheels for
in-container mobility. HydroVolta's public \salinbloc brochure states
container throughput by format: a \textbf{20\,ft} unit handles
2--30\,m\textsuperscript{3}/h ($\approx$44--660\,m\textsuperscript{3}/day
at 22\,h/day operation); a \textbf{40\,ft} unit handles
15--60\,m\textsuperscript{3}/h ($\approx$330--1{,}320\,m\textsuperscript{3}/day)
--- see \cref{sec:fieldperformance} and \cref{ch:appparams}. No container
weight figure was found in any source reviewed.

\hvfig{fig-nameplate.png}{Suggested figure: photograph of the unit's
  nameplate once affixed, showing serial number, manufacture date,
  voltage/power rating, IP rating and CE mark.}{fig:nameplate}


\chapter{Installation}
\label{ch:install}
\chaptertoc

\section{Delivery and unloading}
TODO: no document describing crane/forklift unloading, container weight,
or lifting-point locations was found. The container general-arrangement
drawings (2308CA1/2308CA2, \cref{ch:appdrawings}) likely show lifting
points but have not yet been reviewed for this purpose.

\hvfig{fig-delivery-unloading.png}{Suggested figure: photograph or diagram
  of the container being lifted/unloaded on site, with lifting points
  marked, once the procedure is confirmed.}{fig:delivery}

\section{Site requirements}
The system is built around one or two ISO shipping containers
(pretreatment~+ treatment, \cref{ch:overview}), 20\,ft/40\,ft, High-Cube or
Standard. In HydroVolta's documented BASF pilot, the pretreatment unit used
a 10\,ft container. The site must provide a water source (the pilot drew
from a cooling-tower basin via a submersible pump) and space alongside the
containers for a drain tank. TODO: ambient temperature range, ground/
foundation/leveling requirements and clearances beyond what the pilot
layout implies were not found in a dedicated site-requirements document.

\section{Positioning the container}
TODO: no explicit positioning/orientation document was found. From the
piping documented for commissioning (\cref{sec:watercx}), the pretreatment
container's filtered-water outlet (2") must be within pipe-run distance of
the treatment container's inlet (2"), and both containers' 3" drain
outlets should be positioned so they can be tee'd together to a shared
drain tank.

\section{Water connections}
\label{sec:watercx}
\begin{itemize}
\item \textbf{Pretreatment container:} feed inlet 2\textonehalf" (M) from
  the raw water source (in the pilot, a submersible pump rated
  4.5\,m\textsuperscript{3}/h at 6\,bar); filtered-water outlet 2" (M);
  drain outlet 3" (M).
\item \textbf{Treatment container:} inlet 2" (M), fed from the
  pretreatment filtered-water outlet; product-water outlet 2" (M); drain
  outlet 3" (M).
\item Pump-to-pretreatment feed line: PVCu, 60\,mm, terminated 2\textonehalf" (F).
\item Pretreatment-to-treatment line: PVCu, 40\,mm.
\item Both containers' 3" (M) drain outlets combine via a tee into one
  line to the shared drain tank.
\item Fit O-rings and PTFE tape at every threaded joint.
\end{itemize}
Key in-line pretreatment components: an ultrafiltration unit (AKKIM
APM~80) via 2" Victaulic couplings, and disc filters (AZUD DF100/DF50/DF20)
via 2" (F) 63\,mm transition pieces.

\hvfig{fig-water-connections.png}{Suggested figure: a labeled diagram of
  the pretreatment/treatment container pair showing every hydraulic
  connection point (feed, filtered-water, product, drain) with its size
  and thread type.}{fig:water-connections}

\section{Electrical connection}
The control cabinet is IP55-rated. Incoming supply on the documented pilot
was 3$\times$400\,V\,AC (3P+N+PE, per IEC~60445), max.\ 25\,A, with a 24\,V\,DC
control bus for valves and I/O (from the pretreatment electrical schematic,
sheet~E1/document~I2). A separate electrical schematic series exists
specifically for a \textbf{treatment container} (as opposed to
pretreatment) under the BASF pilot project (folder \texttt{2\_Electrical/
0\_Treatment container}, multiple dated AutoCAD Electrical \texttt{.dsn}
revisions) --- TODO: this series was located but not opened (CAD format);
a human with AutoCAD Electrical should open the latest revision for the
treatment-container-specific breaker/RCD sizing, since the figures above
are from the pretreatment side only. TODO: confirm whether either value
set applies to the standard \salinbloc product or is pilot-specific.

\section{Network and telemetry connection}
See \cref{ch:telemetry} for the router, remote-access gateway and cloud
dashboard used for network/telemetry connectivity.

\section{Installation checklist}
\begin{procedure}
\item Confirm all hydraulic connections listed in \cref{sec:watercx} are
  complete: pretreatment feed/outlet/drain, treatment inlet/outlet/drain,
  combined drain-to-tank tee.
\item Confirm electrical supply matches the cabinet's rated voltage/current
  (\cref{ch:cabinet}) before energizing.
\item Confirm network/telemetry connection is established (\cref{ch:telemetry}).
\end{procedure}
TODO: no standalone installation checklist document was found; the list
above is assembled from the commissioning data sheet and should be
reviewed against a real site installation before publishing.


\chapter{Commissioning}
\label{ch:commissioning}
\chaptertoc

\section{Before the first start}
Confirm every hydraulic connection in \cref{sec:watercx} is complete and
leak-checked before energizing the cabinet. TODO: no document describing a
formal pre-start electrical isolation/lockout check or inspection
checklist was found beyond the piping steps.

\section{Filling and priming}
Install the submersible feed pump in the source water body; connect its
outlet to the pretreatment 2\textonehalf" (F) inlet with 60\,mm PVCu pipe;
connect the pretreatment filtered-water outlet to the treatment container
inlet with 40\,mm PVCu pipe; join both containers' drain outlets via a tee
to the drain tank; verify the product-water outlet is connected per site
requirements. TODO: no document states specific priming steps (venting,
air removal, initial fill sequence) beyond making these connections.

\section{First power-up}
TODO: no formal first-power-up procedure (breaker/contactor closing order,
PLC boot checks, HMI startup) was found. What \emph{does} exist is a real
commissioning engineer's log from a pilot startup, which illustrates the
general shape of a ramp-up sequence (source: pilot commissioning notes,
HydroVolta engineering library) --- treat it as an informative example, not
a validated procedure for the standard product:
\begin{procedure}
\item Start the feed and process pumps at a low VFD frequency
  ($\sim$35\,Hz in the logged example) to establish flow gently rather
  than at full speed.
\item Increase the dosing-pump VFD frequency in steps (e.g.\ to
  $\sim$60\,Hz) once feed flow is confirmed stable.
\item Confirm each stream reaches a stable pressure and flow reading
  before proceeding --- the log explicitly checks \enquote{stable
  pressure} and \enquote{stable flow} at each stack position before
  applying stack current.
\item Apply DC current to the stack(s) only once hydraulics are stable,
  and confirm conductivity in the target stream stabilizes at the
  expected value.
\item Bring additional stacks online one at a time if the skid has more
  than one --- the log notes that current draw on an already-running
  stack visibly drops as a newly started stack begins drawing current,
  which is expected stack-to-stack interaction, not a fault.
\end{procedure}

\section{Setting the operating parameters}
\label{sec:opparams}
Dosing pumps are controlled via a 4--20\,mA analog signal from ADAM I/O
modules. Default configuration: \textbf{SET~1} = 4.0\,mA (minimum, pump
stopped, 0 impulses/min), \textbf{SET~2} = 20.0\,mA (maximum,
180\,impulses/min); internally, a raw integer of 0 corresponds to 3.8\,mA
and 4095 to 19.7\,mA. Each dosing channel is verified across 4/8/12/16/20\,mA.
TODO: hydraulic setpoints (target flow/pressure) and stack-voltage
setpoints are referenced in HydroVolta's internal user-manual outline but
no actual values were found --- see also \cref{ch:appparams}.

\section{Functional test}
A real, filled-in functional test record exists for the BASF pilot's
pretreatment dosing pumps and control panel:
\begin{procedure}
\item Verify each dosing channel's SET~1/SET~2 defaults (4.0\,mA / 0
  impulses/min, 20.0\,mA / 180\,impulses/min) and confirm out-of-range
  signals (below 4\,mA, above 20\,mA) trigger a stop.
\item Test every alarm digital input and dosing digital output for correct
  read/operate behavior.
\item Drive each dosing pump across 4/8/12/16/20\,mA and confirm the
  correct response at each step.
\end{procedure}
TODO: no document describes functional tests for the treatment
(\sonixed/stack) side, pumps, or sensors beyond pretreatment dosing --- see
the relevant Part~II chapters for what should be tested.

\section{Handover}
TODO: no document describing a formal handover procedure, sign-off sheet,
or customer training checklist was found.


\chapter{Daily operation}
\label{ch:operation}
\chaptertoc

\begin{note}
This chapter now draws on HydroVolta's own \emph{HMI Description}
specification (revision~V5, Morocco project, dated 2026-06-16) --- the
design document the automation team is building the PLC program against.
It describes screens, controls and workflow \emph{as specified}; the PLC
program itself was still in development at the time of writing
(\cref{ch:about}), so exact on-screen wording/layout on a specific
delivered unit may differ slightly from what is described here. Where an
older pilot's tag names differ from this specification (e.g.\
\texttt{TSA}-prefixed tags used elsewhere in this manual), both are noted.
\end{note}

\section{Operating elements}
The local control panel exposes physical Start, Stop and Reset controls, a
Manual/local selector, a feedwater-source selector, and a CIP-mode input
(control-system tags \plctag{Start}, \plctag{Stop}, \plctag{Reset},
\plctag{Manual}, \plctag{FW\_Select}, \plctag{LS\_Select}, \plctag{CIP}),
alongside a four-color status LED stack (\plctag{GREEN\_LED},
\plctag{BLUE\_LED}, \plctag{ORANGE\_LED}, \plctag{RED\_LED}). The HMI
specification additionally confirms a \textbf{hardware Emergency Stop
button that bypasses the PLC entirely} --- see \cref{sec:estop}.

\hvfig{fig-local-control-panel.png}{Suggested figure: photograph of the
  local control panel with each button, selector switch and LED labeled
  and cross-referenced to the tag names above.}{fig:local-panel}

\section{Plant modes}
Per the HMI specification, the unit operates in one of nine defined
modes, selected from the \textbf{Mode Control} screen (\cref{sec:hmiscreens}):
\begin{enumerate}
\item \textbf{Initial fill} --- automatic filling of the system before
  first production.
\item \textbf{Production (forward)} --- normal automatic production.
\item \emph{(reversal is a timed sub-state of production, not a separate
  numbered mode --- see \emph{Starting the unit} below)}
\item \textbf{Controlled shutdown} --- an orderly stop sequence.
\item \textbf{CIP} --- automatic clean-in-place, after a CIP-PREP step
  (High~pH/NaOH or Low~pH/HCl, \cref{ch:dosing}) has conditioned the loop.
\item \textbf{Flush/rinse} --- freshwater flush of the process loop.
\item \emph{(reserved)}
\item \textbf{Manual} --- individual device control, access-gated (see
  \emph{Manual mode} below).
\item \emph{(reserved)}
\end{enumerate}
A separate \textbf{Safe Stop} state exists outside the numbered mode list:
it is entered automatically on certain faults and requires the recovery
procedure in \emph{Emergency stop and restart} below before the operator
can select a new mode.

\section{Starting the unit}
Per the HMI specification, starting the unit is a two-step mode selection
from the \textbf{Mode Control} screen: first \textbf{Start AUTO INITIAL
FILL}, then, once initial fill completes, \textbf{Start AUTO PRODUCTION}.
Production runs in \emph{forward} flow with a periodic electrodialysis-
reversal step (\cref{ch:stacks}); the Overview screen shows a countdown
timer to the next reversal. TODO: the specification confirms the button
sequence and mode names but does not itself state numeric pre-start
interlocks/permissives (e.g.\ minimum tank levels) --- those are defined
in the PLC program, still in development.

\section{HMI screens}
\label{sec:hmiscreens}
Per the HMI specification, the operator interface is organized into six
screens:
\begin{itemize}
\item \textbf{Overview} --- current plant mode and Safe Stop indicator;
  key performance indicators (product conductivity, product flow,
  recovery); product routing (to customer vs.\ recirculation); a status
  badge for each stack (power on/off, current, conductivity,
  pressure-OK); pump run/fault status; the position of the critical
  reversal and CIP-loop valves; the reversal countdown timer; and the
  active alarm count with highest severity.
\item \textbf{Tank Summary} --- level and conductivity for every tank,
  and the state of every fill/drain valve.
\item \textbf{Mode Control} --- the plant-mode buttons described above,
  plus \emph{Drain Concentrate Tank}, \emph{Switch to Manual}, \emph{ACK
  Alarm} and \emph{Reset Trip}.
\item \textbf{Alarm Summary} --- active, acknowledged and historical
  alarms, each showing its code, severity, cause, effect, timestamp and a
  recommended action, following the ISA-18.2 alarm-management convention
  --- see \emph{Alarm handling} in \cref{ch:maintenance}.
\item \textbf{Maintenance / Manual} --- individual device commands
  (pumps, valves, dosing pumps, stack power), gated by the access level
  in \emph{Manual mode} below.
\item \textbf{Trends} --- historical KPI graphs.
\end{itemize}

\hvfig{fig-hmi-home-screen.png}{Suggested figure: screenshot of the HMI
  Overview screen, once the PLC program is complete and running on an
  actual unit, annotated with each element's function per the list
  above.}{fig:hmi-home}

\section{Reading the process values}
Confirmed process readings available to the operator include pH,
pressure, flow and conductivity/level for the sensor array (\plctag{TSA\_pH},
\plctag{TSA\_pressure}, \plctag{TSA\_Flow}, \plctag{TSA\_Cond},
\plctag{TSA\_Level}), plus conductivity and level for the~TB tank
(\plctag{TB\_Cond}, \plctag{TB\_Level}) and level for the~TE (electrolyte)
tank (\plctag{TE\_Level}) --- see \cref{ch:tanks,ch:sensors}. Additional
electrical metering is available via power-meter tags (\plctag{PM3150}).

\begin{note}
HydroVolta's I/O tag lists follow a readable naming convention that helps
when reading raw tag names on an HMI or historian trend: a leading letter
indicates the measurement type (\texttt{P}~= pressure, \texttt{FL}~=
flow), followed by \texttt{c}/\texttt{d}/\texttt{e} for
concentrate/diluate/electrolyte, and \texttt{S1}--\texttt{S4} for the
stack number when the tag is per-stack. Tank tags use \texttt{T} plus a
capital letter for the tank's role (\texttt{TD}~= diluate/feed tank,
\texttt{TC}~= concentrate loop tank, \texttt{TE}~= electrolyte tank,
\texttt{TB}~= product/breaking tank). Valve tags use \texttt{VA} plus a
functional group number. TODO: confirm whether the standard \salinbloc
product uses this exact convention or the \texttt{TSA}-prefixed scheme
seen elsewhere in this manual --- both appear in HydroVolta's pilot
projects, and the finished PLC program (still in development, see
\cref{ch:about}) will determine which applies to a given unit.
\end{note}

\section{Changing setpoints}
Per the HMI specification, every setpoint, threshold, timer, dose and
recipe value is adjustable directly from the HMI at runtime, without
re-downloading the PLC program. Which parameters a given user can change
is governed by a three-level access hierarchy:
\begin{itemize}
\item \textbf{Operator:} monitor screens, acknowledge alarms, start/stop
  production, drain the concentrate tank, run a manual flush, view
  trends. Cannot change setpoints.
\item \textbf{Maintenance:} everything an Operator can do, plus
  individual pump/valve/dosing-pump control, manual fills/drains, and
  non-recipe parameter changes.
\item \textbf{Admin/Engineer:} everything Maintenance can do, plus
  stack-power control (as a set, not individually), recipe parameters,
  cloud command flags, and user-account management.
\end{itemize}
TODO: the specific PIN/login mechanism for each access level, and the
numeric default/range for every adjustable parameter, are defined in the
full HMI specification's parameter table but were not transcribed into
this manual --- consult HydroVolta engineering for the complete parameter
table if commissioning a unit before this manual's next revision.

\section{Manual mode}
A \plctag{Manual} digital input and \plctag{LS\_Select}/\plctag{FW\_Select}
mode-selection inputs exist in the control I/O map, and the HMI
specification confirms \textbf{Manual} as Mode~8, reached via \emph{Switch
to Manual} on the Mode Control screen and gated to the Maintenance/Admin
access levels above --- an Operator-level user cannot enter Manual mode.
In Manual mode, the Maintenance/Manual screen exposes direct command of
every process pump, every 2-way and 3-way valve, all five dosing pumps,
and stack power (as a set only, even in Manual). One routing valve
(\plctag{VAR.S.5} in the specification's tag scheme) is explicitly
\textbf{AUTO-only} and cannot be operated from Manual mode under any
access level --- do not attempt to command it manually.

\section{Stopping the unit}
A \plctag{Stop} digital input exists on the local control panel, and the
HMI specification's \textbf{Controlled Shutdown} (Mode~4) is the normal,
orderly way to stop production from the Mode Control screen --- prefer it
over simply removing power, so the system can close valves and stop pumps
in a defined sequence rather than an uncontrolled one. TODO: the exact
step-by-step shutdown sequence executed by Mode~4 is internal to the PLC
program and was not itself transcribed into this manual.

\section{Emergency stop and restart}
\label{sec:estop}
The HMI specification confirms a \textbf{hardware Emergency Stop button
that bypasses the PLC entirely} --- pressing it removes power/motion
directly, independent of software state. TODO: its physical location on
the control panel/cabinet was not confirmed in the documents reviewed;
verify on the actual unit and update \cref{fig:local-panel} accordingly.
Certain faults also trip the system into a software \textbf{Safe Stop}
state (distinct from the hardware E-stop) --- both require the following
recovery sequence before production can resume, per the HMI specification:
\begin{procedure}
\item Investigate the cause using the Alarm Summary screen
  (\cref{sec:hmiscreens}) --- each entry shows its cause, effect and a
  recommended action.
\item Physically inspect and resolve the underlying condition.
\item Press \textbf{ACK Alarm} on the Mode Control screen.
\item Press \textbf{Reset Trip} --- this button is only enabled once the
  fault condition has actually cleared.
\item Select the next mode (e.g.\ Production). The PLC re-validates all
  permissives before allowing the mode to start, and will reject the
  request with a message if a required condition is still not met.
\end{procedure}
\begin{danger}
Never attempt to defeat, bridge, or otherwise bypass the hardware
Emergency Stop button. If the E-stop has been pressed, do not restart the
unit until the cause has been fully investigated and resolved.
\end{danger}


\chapter{Maintenance and troubleshooting overview}
\label{ch:maintenance}
\chaptersubtitle{General schedule and fault-finding approach --- for a
  specific compartment's checks, specifications and fault table, see its own
  chapter in Part~II.}
\chaptertoc

\section{Safety during maintenance}
Before performing any wiring or sensor checks, isolate and verify the
control circuit: control voltage is referenced to ground at 24\,V\,DC ---
confirm 24\,V\,DC is present between the terminal and PE before assuming a
wiring fault, and never connect a sensor to the PLC before bench-testing it
(\cref{ch:sensors}). TODO: no dedicated safety document (lockout/tagout,
PPE, chemical/electrode handling) specific to \salinbloc maintenance was
found --- see \cref{ch:safety}.

\section{Maintenance schedule}
No complete maintenance interval table (daily/weekly/monthly/annual) was
found for the \salinbloc system as a whole. The one concrete, recurring
interval documented anywhere in HydroVolta's engineering library is the
CIP (clean-in-place) cycle described in \cref{ch:dosing}, which ran
approximately \textbf{weekly} in the documented pilot; the pretreatment
UF stage is otherwise cleaned on a TMP-triggered basis rather than a fixed
interval (\cref{ch:pretreatment}).

\begin{notice}
The table below is a \textbf{recommended starting schedule}, not a
HydroVolta-issued one. The two rows marked \enquote{Confirmed} come
directly from source documents cited elsewhere in this manual; every
other row is this manual's synthesis of general good practice for
equipment of this type (visual/level checks, periodic sensor
verification, periodic actuator exercise, annual inspection), pending
review and correction by HydroVolta engineering once real field data
exists. Do not treat the unconfirmed rows as validated intervals.
\end{notice}

\hvtablebegin{L C{20mm} L}{Recommended starting maintenance schedule}{tab:maintschedule}
\hvhead{\textbf{Task} & \textbf{Interval} & \textbf{Source / basis}}\\\midrule
Visual walk-round: check for leaks, unusual noise/vibration, tank levels, HMI alarm state & Daily & Recommended practice (unconfirmed)\\
Record process values in maintenance log (\cref{ch:appparams} ranges) & Daily & Recommended practice (unconfirmed)\\
CIP cycle (acid + hypochlorite steps, \cref{ch:dosing}) & Weekly & \textbf{Confirmed} --- documented pilot interval\\
Check dosing chemical tank levels and top up & Weekly & Recommended practice (unconfirmed)\\
UF chemically enhanced backwash (CEB) & When TMP rises $\geq$0.9\,bar above initial & \textbf{Confirmed} --- UF manufacturer spec (\cref{ch:pretreatment})\\
Bench-verify each 4--20\,mA sensor loop (\cref{sec:sensorcal}) & Quarterly, or on any suspect reading & Recommended practice (unconfirmed)\\
Exercise/inspect valve actuators (manual override, seat wear) & Quarterly & Recommended practice (unconfirmed)\\
Inspect control cabinet (dust ingress on IP55 enclosure, terminal torque) & Semi-annually & Recommended practice (unconfirmed)\\
Full system inspection and log review & Annually & Recommended practice (unconfirmed)\\
\bottomrule\end{xltabular}\par
TODO: every compartment's \emph{Maintenance} section in Part~II should be
reviewed by HydroVolta engineering and this table corrected/expanded with
real intervals as field data accumulates.

\section{Maintenance log}
A maintenance log should record, per event: date, technician,
equipment/asset~ID (e.g.\ stack, pump, sensor tag), the check or service
performed, readings taken (pressure, flow, pH, conductivity, etc.) versus
expected range, parts replaced (with part number), and any follow-up
action required. Its purpose is a traceable per-unit service history so
recurring faults, consumable replacement intervals and warranty-relevant
events can be reviewed over the system's lifetime.

\hvtablebegin{L C{22mm} C{30mm} L C{20mm}}{Maintenance log template}{tab:maintlog}
\hvhead{\textbf{Date} & \textbf{Technician} & \textbf{Equipment / tag} & \textbf{Check / service performed} & \textbf{Result}}\\\midrule
 & & & & \\
 & & & & \\
\bottomrule\end{xltabular}\par

\section{Alarm handling}
\label{sec:alarmhandling}
Per HydroVolta's HMI design specification (\cref{ch:operation}), alarms
follow the ISA-18.2 alarm-management convention and always resolve the
same way, regardless of which alarm fired:
\begin{procedure}
\item The PLC raises the alarm; it appears on the \textbf{Alarm Summary}
  screen with its code, severity, cause, effect, timestamp and a
  recommended action.
\item The operator investigates using that information and this manual's
  Troubleshooting section for the relevant compartment (below).
\item Once the operator has acknowledged they are aware of it, press
  \textbf{ACK Alarm} (Mode Control screen) --- the alarm moves to the
  acknowledged list.
\item Physically resolve the underlying condition.
\item Press \textbf{Reset Trip} --- only enabled once the condition has
  actually cleared, so a reset that won't take is itself informative: it
  means the underlying condition has not actually gone away yet.
\item Resolved alarms remain in the alarm \emph{history} with their full
  timestamp trail, for the maintenance log (\cref{ch:maintenance}).
\end{procedure}
A fault severe enough to trip \textbf{Safe Stop} needs the fuller recovery
sequence in \cref{sec:estop}, not just ACK~+ Reset.

\section{How to use the troubleshooting sections}
Each Part~II chapter ends with a \emph{Troubleshooting} section scoped to
that compartment: symptom, likely cause, and remedy. Start with the
compartment most closely tied to the symptom (e.g.\ a conductivity alarm
$\rightarrow$ \cref{ch:sensors} or \cref{ch:tanks}; a communication timeout
$\rightarrow$ \cref{ch:cabinet} or \cref{ch:telemetry}). If the fault
spans several compartments or persists after following the relevant
section, escalate per \cref{sec:whentocallservice}.

\section{When to call service}
\label{sec:whentocallservice}
Escalate to HydroVolta service when: a sensor fails its bench test
(\cref{ch:sensors}) and no spare is available on site; a remote-access
device (e.g.\ the ADAM-3600 RTU) remains unreachable after the documented
recovery steps (\cref{ch:cabinet}) and is not even pingable; or process
readings remain outside the documented operating envelope after checking
feedwater quality, pump duty point and pipe sizing. TODO: no source
document defines formal escalation criteria, response-time commitments, or
service contact information --- see also \cref{ch:about} (Contact and
support).


%% ---- Part II: Container compartments -------------------------------------
%% Every chapter in this part follows the same five-section shape so a
%% reader (or a technician printing just this chapter) gets a complete,
%% self-contained reference for that one piece of equipment:
%%   Overview and function / Key specifications / Operation /
%%   Maintenance / Troubleshooting

\chapter{Pretreatment and Filtration Skid}
\label{ch:pretreatment}
\chaptersubtitle{Disc filters, ultrafiltration unit and pretreatment piping}
\chaptertoc

\section{Overview and function}
The pretreatment container conditions raw feed water before it reaches the
\sonixed stacks: disc filtration removes coarse suspended solids, and an
ultrafiltration (UF) stage removes fine particulates and reduces fouling
potential. Filtered water passes to the treatment container
(\cref{ch:overview}); rejected solids and backwash water are routed to the
drain tank.

\hvfig{fig-pretreatment-skid.png}{Suggested figure: photograph of the
  assembled pretreatment skid inside the container, with the disc filters,
  UF unit and main piping runs labeled.}{fig:pretreatment-skid}

\section{Key specifications}
\begin{itemize}
\item \textbf{Ultrafiltration:} Akualys/AKKIM APM~80 hollow-fiber UF
  module(s), 2" Victaulic connections. Two modules are confirmed ordered
  for HydroVolta (160\,m\textsuperscript{2} total membrane area). Membrane:
  modified PVDF hollow fiber, outside-to-inside flow, fiber OD/ID
  1.4/0.8\,mm, nominal MWCO $\leq$150{,}000\,Da ($\approx$30\,nm pore
  size), 80\,m\textsuperscript{2} active area per module.
\item \textbf{UF feedwater limits:} temperature $\leq$40\,\textdegree C
  (typ.\ 25\,\textdegree C); turbidity $\leq$250\,NTU (typ.\ 50\,NTU); TSS
  $\leq$80\,mg/l; TOC $\leq$30\,mg/l; pH 6--9; instantaneous free chlorine
  $\leq$150\,mg/l.
\item \textbf{UF operation:} dead-end/crossflow, filtrate flux
  45--180\,l/m\textsuperscript{2}h at 25\,\textdegree C, flow capacity
  3.60--14.40\,m\textsuperscript{3}/h per module, feed inlet pressure
  2--3\,bar (instantaneous max.\ 5\,bar), transmembrane pressure (TMP)
  0.4--2\,bar, filtrate SDI $\leq$2.5, filtrate turbidity $\leq$0.2\,NTU.
\item \textbf{Disc filtration:} AZUD Helix Automatic self-cleaning disc
  filters (FT-series, automatic backwash), 100-micron mesh discs; the
  documented BASF pilot P\&ID additionally references disc-filter
  positions DF100/DF50/DF20 by micron rating --- TODO confirm which AZUD
  Helix model number(s) map to the standard \salinbloc pretreatment train.
\item Feed inlet: 2\textonehalf" (M); filtered-water outlet: 2" (M); drain:
  3" (M).
\item Feed pump (pilot configuration): submersible, 4.5\,m\textsuperscript{3}/h
  at 6\,bar.
\end{itemize}
TODO: confirm whether the AKKIM/AZUD models above are the standard
\salinbloc bill-of-materials or pilot-specific; a finalized
pretreatment BOM was not found. Several other filtration technologies
(hydrocyclone separator, automatic screen filters, high-efficiency sand
filters, microfilters) exist as vendor catalogs in HydroVolta's component
library but are not confirmed as installed on any \sonixed pretreatment
train.

\section{Operation}
No dedicated pretreatment-specific operating procedure beyond the general
startup/shutdown described in \cref{ch:operation} was found, but the UF
module manufacturer documents the following normal operating cycle, which
is the closest available reference for the standard \salinbloc pretreatment
skid:
\begin{procedure}
\item \textbf{Production:} a brief forward-flush step, then filtration
  (dead-end or crossflow) until the trigger for backwash is reached.
\item \textbf{Backwash:} drain, then air-scour (max.\
  20\,Nm\textsuperscript{3}/h at 600--800\,mbar) to loosen the fouling
  layer, then backwash bottom-to-top and top-to-bottom
  (backwash pressure max.\ 2.5\,bar) to flush it out.
\end{procedure}
TODO: confirm the actual backwash trigger used on \salinbloc (pressure
differential vs.\ timer) and the cycle duration/frequency; the manufacturer
spec above describes the module's capability, not HydroVolta's configured
setpoint.

\section{Maintenance}
The UF manufacturer documents a Chemically Enhanced Backwash (CEB) /
CIP cycle, triggered when TMP rises $\geq$0.9\,bar above its initial
value:
\begin{procedure}
\item Run an extended backwash cycle (as in \emph{Operation} above), then
  soak the module in cleaning chemical for part of the cycle.
\item CEB chemical concentrations: NaOCl 1000\,mg/l, NaOH 500\,mg/l, HCl
  1000\,mg/l, or citric acid 1--2\% --- select per the fouling type
  (organic/biological vs.\ scale) per \cref{ch:dosing}.
\item Total CIP duration $\approx$2\,hours, at a cleaning flow of
  1--2\,m\textsuperscript{3}/h per module.
\end{procedure}
For extended shutdowns (\cref{ch:shutdown}), the manufacturer specifies a
glycerol/sodium-metabisulfite preservation soak, and requires a minimum
4-hour flush with $\geq$14\,m\textsuperscript{3} of feed water before
restart. Before returning a preserved or newly installed module to
service, disinfect with 100\,ppm NaOCl circulated for 10\,minutes, held
for 1\,hour, then drained until the filtrate runs chlorine-free. TODO: no
document gives a routine cleaning-interval schedule (as opposed to the
TMP-triggered CEB above) --- see \cref{ch:maintenance} for the general log
template.

\hvfig{fig-filter-cleaning.png}{Suggested figure: step-by-step photos of
  the disc-filter backwash and UF CEB/CIP procedure above, once
  photographed on an actual unit.}{fig:filter-cleaning}

\section{Troubleshooting}
\begin{itemize}
\item \textbf{Low filtered-water flow / high pressure drop:} check pump
  duty point against the 4.5\,m\textsuperscript{3}/h @ 6\,bar rating and
  pipe sizing (60\,mm feed line, 40\,mm filtered-water line) before
  suspecting the filters themselves.
\item \textbf{Rising transmembrane pressure (TMP):} if TMP has risen
  $\geq$0.9\,bar above its initial value, run the CEB/CIP cycle above
  rather than waiting for a scheduled cleaning.
\item \textbf{Poor downstream water quality:} verify feed water against
  the operating envelope in \cref{ch:stacks} (turbidity, iron, SDI) and
  against the UF feedwater limits above --- an out-of-envelope feed is a
  likely root cause on either stage.
\end{itemize}
TODO: no documented fault table (e.g.\ specific alarm codes) for the disc
filters or UF unit was found beyond the operating-envelope figures above.


\chapter{Feed and Process Pumps}
\label{ch:pumps}
\chaptersubtitle{Feed, diluate and concentrate pumps on each SonixED skid}
\chaptertoc

\section{Overview and function}
Each \sonixed skid carries three process pumps --- feed, diluate and
concentrate --- that circulate water through the stack(s) at the required
flow and pressure. Pumps are magnetic-driven (sealless) and inverter-
(VFD-)controlled so flow can be adjusted without throttling valves.

\hvfig{fig-process-pumps.png}{Suggested figure: photograph of the three
  process pumps on a skid, labeled feed/diluate/concentrate, with their
  VFD units visible.}{fig:process-pumps}

\section{Key specifications}
\begin{itemize}
\item Always 3 pumps per skid (feed, diluate, concentrate).
\item Magnetic-driven, inverter-driven, polypropylene/stainless-steel
  construction, seawater-compatible.
\item Pretreatment feed pump (pilot): Grundfos CRN~10-5 A-FGJ-A-E-HQQE,
  flanged connections.
\item Rated hydraulic flow per stack: 24\,m\textsuperscript{3}/day
  (SonixED-S datasheet).
\item Each pump is driven by its own VFD (control tags
  \plctag{VFD\_1\_D1}\dots\plctag{VFD\_6\_E2} in the pilot configuration).
\end{itemize}

\section{Operation}
Pump speed is set by the PLC via the VFDs to hold the process at its
target flow/pressure; TODO no document confirms the exact control loop
(open-loop speed vs.\ closed-loop flow/pressure control). Manual
speed override, if available, is reached via the \plctag{Manual} mode
described in \cref{ch:operation}.

\section{Maintenance}
TODO: no pump-specific maintenance schedule (bearing/seal inspection,
impeller check interval) was found. Magnetic-driven pumps have no shaft
seal to replace, which reduces routine maintenance versus a conventional
centrifugal pump, but confirm this against the actual pump model fitted.

\section{Troubleshooting}
Each pump's VFD reports a \plctag{FaultWord} that should be checked first
on any pump-related alarm. Typical causes to check, in order:
\begin{itemize}
\item \textbf{Pump not running / VFD fault:} check \plctag{FaultWord} on
  the relevant VFD tag; check incoming supply and control-voltage presence
  (\cref{ch:cabinet}).
\item \textbf{Low flow / high pressure:} check for blockage upstream
  (\cref{ch:pretreatment}) or downstream (stack fouling,
  \cref{ch:stacks}) before assuming a pump fault.
\item \textbf{Pump running but no output:} check for air lock (magnetic-
  driven pumps are not self-priming) --- TODO confirm priming procedure.
\end{itemize}


\chapter{SonixED Stack Modules}
\label{ch:stacks}
\chaptersubtitle{The electrodialysis stacks at the heart of the unit}
\chaptertoc

\section{Overview and function}
The \sonixed stack is the unit's core: an assembly of alternating
cation- and anion-exchange membranes, titanium electrodes and spacers,
with integrated ultrasound transducers for chemical-free antifouling.
Under an applied DC voltage, ions migrate across the membranes from the
diluate stream into the concentrate stream, desalinating the diluate. Each
skid carries 1--3 stacks (standard configurations 1SS/2SS/3SS).

\hvfig{fig-stack-module.png}{Suggested figure: labeled photograph of a
  single \sonixed stack module, showing the electrode housings, membrane
  stack, ultrasound transducers and 1" BSP connections.}{fig:stack-module}

\hvfig{fig-stack-cutaway.png}{Suggested figure: a cutaway/exploded diagram
  of one cell pair, showing the diluate/concentrate/electrolyte flow paths
  through alternating membranes --- useful for explaining the process to a
  new operator.}{fig:stack-cutaway}

\section{Key specifications}
\begin{itemize}
\item Module dimensions (l$\times$b$\times$h): 300\,$\times$\,400\,$\times$\,1060\,mm;
  weight $\sim$150\,kg; connection 1" BSP.
\item Flowrate per stack: 0.5--10\,m\textsuperscript{3}/h (skid datasheet);
  24\,m\textsuperscript{3}/day (stack datasheet).
\item Salt rejection: up to 90\%; water recovery: up to 98\%.
\item Operating pressure: 0.5--2.5\,bar (typical 0.5--1\,bar); maximum
  3\,bar.
\item Temperature: 5--40\,\textdegree C (stack datasheet) / 5--55\,\textdegree C
  (skid datasheet) --- TODO reconcile which applies to the finished
  product.
\item DC voltage across stack: 60--600\,V; DC current: 3--20\,A.
\item Cell pairs per stack: 50--250; effective membrane area: 10--50\,m\textsuperscript{2}
  (\enquote{long} variant) or 20--100\,m\textsuperscript{2} (\enquote{square} variant).
\item Membrane: Fujifilm Type~10, polyolefin.
\item Feed water requirements: salinity $\leq$90{,}000\,ppm; pH 5--8.5
  typical (tolerable 2--13); turbidity $<$1\,NTU; suspended solids
  $<$20\,\textmu m; SDI(5\,min) 10; iron $<$0.5\,mg/l; free chlorine
  0.5\,ppm continuous / 30\,ppm intermittent tolerance.
\end{itemize}

\section{Operation}
Stacks operate continuously under PLC control once commissioned
(\cref{ch:commissioning}); electrodialysis \emph{reversal} periodically
swaps stream polarity as part of the ultrasound antifouling strategy ---
TODO confirm the reversal interval/trigger. Process values (pressure, flow,
conductivity per stream) are monitored via the sensors described in
\cref{ch:sensors}.

\section{Maintenance}
The membrane is the stack's principal wear item (weight $\sim$150\,kg,
dimensions and connection as in \emph{Key specifications} above). TODO:
part number, reorder code and recommended replacement interval were not
found; no CIP or cleaning procedure specific to the \sonixed stack was
found among the documents reviewed.

\section{Troubleshooting}
If salt rejection or recovery drops noticeably below the rated ranges
above, check feedwater against the operating envelope in \emph{Key
specifications} first --- turbidity, iron and free chlorine excursions are
the most likely root causes of fouling-related performance loss. TODO: no
documented fault table linking specific water-quality deviations to
specific symptoms (e.g.\ \enquote{high iron causes X}) was found beyond the
acceptable operating envelope itself.


\chapter{Valves and Manifold}
\label{ch:valves}
\chaptersubtitle{Automated 3-way valves and the feed distribution manifold}
\chaptertoc

\section{Overview and function}
Automated 3-way valves direct flow between feed, product and recirculation
paths (e.g.\ for electrodialysis reversal, \cref{ch:stacks}), while the
manifold assembly distributes feed water evenly to every stack on a skid.

\hvfig{fig-valve-manifold.png}{Suggested figure: photograph of the
  manifold assembly with its 3-way valves, labeled by function (feed,
  product, recirculation).}{fig:valve-manifold}

\section{Key specifications}
\begin{itemize}
\item 3-way automated valves: PVC construction.
\item Manifold: DN25--DN63, PVC and stainless-steel fittings.
\end{itemize}
Two actuated-valve product lines are confirmed present in HydroVolta's
component library as candidates for this duty (TODO: confirm which is
actually fitted on the standard \salinbloc BOM):
\begin{itemize}
\item \textbf{Georg Fischer (+GF+) e-DIASTAR Type~514} actuated ball
  valve: PVC-U body, 24\,V\,AC/DC multiturn actuator with integrated
  open/closed position feedback, built-in heating element and manual
  emergency override. Available in DN25 (1", PN10) and DN50 (2", PN6)
  sizes matching this manual's connection sizes elsewhere; seal options
  EPDM or PTFE/EPDM.
\item \textbf{B\"urkert Type~3003} electric rotary actuator, fitted to
  quarter-turn ball/butterfly valves: available as on/off or proportional
  control (accepting/outputting 4--20\,mA, 0--20\,mA or 0--10\,V);
  standard manual override, adjustable limit switches, integrated heating
  resistor and torque limiter; UL~94~V0 flame-resistant housing.
\end{itemize}
Both product lines run on the same 24\,V\,DC control bus documented for
the cabinet (\cref{ch:cabinet}), consistent with the troubleshooting
guidance below.

\section{Operation}
Valves are switched automatically by the PLC as part of the stream-
reversal and CIP sequences (\cref{ch:stacks,ch:operation}). TODO: no
document describes manual valve override procedure or fail-safe position
on power loss --- both candidate actuator families above include a manual
override lever, but the fail-safe (last-position vs.\ defined) behavior on
power loss has not been confirmed.

\section{Maintenance}
TODO: no valve-specific maintenance schedule (actuator lubrication, seat
wear inspection interval) was found. See \cref{ch:maintenance} for the
general schedule and log template.

\section{Troubleshooting}
\begin{itemize}
\item \textbf{Valve not switching:} check 24\,V\,DC control voltage at the
  valve terminal against PE before assuming a valve fault
  (\cref{ch:cabinet}).
\item \textbf{Flow imbalance across stacks:} check manifold for blockage
  or a partially closed valve before suspecting a pump or stack issue.
\end{itemize}


\chapter{Dosing System}
\label{ch:dosing}
\chaptersubtitle{Antiscalant, pH correction and CIP dosing pumps}
\chaptertoc

\section{Overview and function}
Dosing pumps inject antiscalant and pH-correction chemicals (NaOH, HCl,
and in the pilot configuration also H\textsubscript{2}SO\textsubscript{4})
into the feed stream, and supply CIP (clean-in-place) chemicals for
periodic cleaning. The documented pilot configuration used 5 dosing
channels.

\hvfig{fig-dosing-skid.png}{Suggested figure: photograph of the dosing
  skid, with each chemical tank and pump labeled (antiscalant, NaOH, HCl,
  CIP).}{fig:dosing-skid}

\section{Key specifications}
\begin{itemize}
\item 5 dosing channels (pilot configuration), each controlled via a
  4--20\,mA signal from ADAM I/O modules.
\item Signal range: SET~1~=~4.0\,mA (0 impulses/min, pump stopped),
  SET~2~=~20.0\,mA (180\,impulses/min).
\item Chemicals confirmed in the pilot tag list and in Safety Data Sheets
  held by HydroVolta: antiscalant (AS), NaOH, HCl, NaOCl (sodium
  hypochlorite) and, in some configurations, H\textsubscript{2}SO\textsubscript{4}
  as an alternative acid.
\item Candidate dosing pump: \textbf{Etatron GB eOne} series --- flow rate
  1--30\,l/h, max.\ pressure 20\,bar, 100--250\,V\,AC 50/60\,Hz supply, max.\
  300 pulses/min, PVDF pump head with double ceramic ball valves and PTFE
  diaphragm. The \textbf{MF} (multifunction) variant adds proportional
  dosing, 4--20\,mA control, level-switch input and failure alarms; the
  \textbf{PLUS} variant integrates a pH/ORP/Cl controller. TODO: confirm
  which variant, and which alternative (Lutz-Jesco MEMDOS Smart /
  EasyChlorGen also appear in HydroVolta's component library), is fitted
  to the standard \salinbloc product.
\end{itemize}

\section{Operation}
See \cref{sec:opparams} in \cref{ch:commissioning} for the default 4--20\,mA
signal configuration and functional-test procedure used to verify each
dosing channel. Setting a dosing pump's frequency from a 4--20\,mA signal
uses two settings: the \textbf{pulse amplitude} (mL per stroke, set
manually at the pump) and the \textbf{pulse frequency} (set remotely as a
\% of the pump's maximum, via the 4--20\,mA signal). Worked example from a
real HydroVolta dosing calculation (antiscalant, 10\,ppm target): dilution
factor 1333, pulse amplitude 0.01--0.1\,mL, giving a dose rate of
$\approx$0.167\,mL/min (36\,mL/h).

\begin{warning}
\textbf{Never dose NaOH into feed or process water} --- HydroVolta's own
CIP procedure documents this explicitly as a scaling risk. NaOH is only
used against pure (permeate-quality) water during CIP.
\end{warning}

A documented weekly CIP cycle (source: HydroVolta pilot CIP procedure)
runs as follows:
\begin{procedure}
\item \textbf{Acid step:} dose HCl to pH~2.5 in the CIP loop, circulating
  30\,minutes forward-flow then 30\,minutes reverse-flow through the
  stack, using pretreated feedwater.
\item Drain the CIP loop to the concentrate tank, then flush with fresh
  feed water for 15--20\,minutes.
\item \textbf{Hypochlorite step:} dose NaOCl to 2\,ppm (effective only
  near pH~7--7.5 --- check/correct pH before this step), circulating
  30\,minutes forward then 30\,minutes reverse.
\item Drain 5\,minutes, then flush with fresh feed water for
  30\,minutes.
\item A variant of this cycle that also treats the electrolyte
  compartment uses NaOH instead of/alongside HCl --- \textbf{only ever
  with pure water}, per the warning above.
\end{procedure}
Typical HCl/NaOH CIP dose rates from a real pilot's dosing calculation
sheet (pure, low-alkalinity water): HCl $\approx$5.5\,mL/min at 1000\,l/h
CIP flow, pulse amplitude 0.1\,mL, $\approx$31\% of max.\ pump frequency;
NaOH $\approx$17.4\,mL/min, pulse amplitude 0.3\,mL, $\approx$32\% of
max.\ frequency. For natural (higher-alkalinity) feed water, acid demand
scales with the water's total alkalinity (TA) --- a seawater-range feed
(TA~25--30\textdegree f) needs roughly 3--4$\times$ the acid dose of a
low-alkalinity feed. TODO: confirm the CIP interval and dose rates that
apply to the standard \salinbloc product; the figures above are a real,
sourced pilot example, not a general specification.

\section{Maintenance}
TODO: no chemical replenishment interval, tubing/diaphragm replacement
schedule, or dosing-pump calibration procedure specific to the standard
product was found. The CIP cycle above ran approximately weekly in the
documented pilot --- use this as a starting reference for the maintenance
schedule in \cref{ch:maintenance} until a formal interval is confirmed.
Per the chemicals' Safety Data Sheets, store all dosing chemicals in their
original, tightly closed containers, away from food/drink, at a
recommended \textbf{15--25\,\textdegree C}; keep incompatible chemicals
(acids vs.\ NaOH vs.\ NaOCl) in separate secondary containment. See
\cref{sec:firstaid} for first-aid and spill-response guidance.

\section{Troubleshooting}
Low-level alarms exist per chemical: \plctag{LL\_Alarm\_AS} (antiscalant),
\plctag{LL\_Alarm\_NaOH}, \plctag{LL\_Alarm\_HCl}, plus per-channel
\plctag{Dosing1\_Alarm}\dots\plctag{Dosing5\_Alarm}. On any dosing alarm,
check chemical tank level first, then verify the 4--20\,mA signal at the
pump terminal using the bench-test method in \cref{ch:sensors}. If a CIP
cycle fails to reach its target pH (2.5 for the acid step), verify the
dosing pump's pulse amplitude/frequency setting against the worked
example above before suspecting a pump fault.


\chapter{Instrumentation and Sensors}
\label{ch:sensors}
\chaptersubtitle{Pressure, flow, pH and conductivity measurement}
\chaptertoc

\section{Overview and function}
Inline sensors give the PLC (and the operator, via the HMI) a continuous
read of the process: pressure, flow, pH and conductivity/level at key
points in the pretreatment and treatment trains, plus a pressure gauge on
each stack. All sensors use a standard 2-wire, loop-powered 4--20\,mA
output.

\hvfig{fig-sensor-array.png}{Suggested figure: photograph of the inline
  sensor array on a skid, with each sensor type (pressure, flow, pH,
  conductivity) labeled.}{fig:sensor-array}

\section{Key specifications}
\begin{itemize}
\item Signal type: 2-wire, loop-powered, 4--20\,mA.
\item Confirmed sensor types in the pilot tag list: pressure
  (\plctag{TSA\_pressure}), flow (\plctag{TSA\_Flow}), pH
  (\plctag{TSA\_pH}), conductivity/level (\plctag{TSA\_Cond},
  \plctag{TSA\_Level}, \plctag{TB\_Cond}, \plctag{TB\_Level},
  \plctag{TE\_Level}).
\item \textbf{Pressure:} WIKA A-10 pressure transmitter (gauge pressure,
  4\,mA~=~0\,bar; 2-wire, IP65/IP67, wetted parts 316L stainless steel,
  non-repeatability $\leq\pm$0.2\% of span). Stated maintenance-free by
  the manufacturer; clean only with water/mild detergent and a soft cloth
  --- repairs are manufacturer-only.
\item \textbf{Flow:} a real, dated quotation (KROHNE Belgium to
  HydroVolta, June~2025, explicitly for \enquote{meetinstrumenten
  onzilting} --- \enquote{desalination measuring instruments}) specifies
  six \textbf{KROHNE OPTIFLUX~1000/1050} electromagnetic flowmeters
  (DN10--150, compact model, max.\ operating pressure 16\,bar, PFA
  lining, Hastelloy~C22 electrodes, IP66/67, product temperature
  $-$25 to +120\,\textdegree C) with an \textbf{IFC~050} flow converter.
\item \textbf{Conductivity:} the same 2025 KROHNE quote specifies
  \textbf{SMARTPAT COND} conductivity sensors.
\item \textbf{pH/ORP:} the same quote specifies \textbf{KROHNE SMARTPAT
  pH/ORP} sensors --- superseding the GF~DryLoc~2751 candidate noted in
  earlier pilots.
\item \textbf{Level:} the same quote specifies an \textbf{80\,GHz
  non-contact radar (FMCW)} level sensor (brand consistent with the
  Holykell/MDS-AM15 candidates found elsewhere, now confirmed as part of
  a coherent, single-vendor 2025 instrumentation package rather than
  scattered evaluation quotes).
\end{itemize}
Because this KROHNE package is dated 2025, covers flow/conductivity/pH/
level in one coherent quotation, and is explicitly scoped to desalination
instrumentation (as opposed to the older, scattered
GF/B\"urkert/Holykell candidates found in earlier pilot projects), it is
the \textbf{best available evidence for the standard \salinbloc
instrumentation package} --- but it is still a quotation, not a
confirmed purchase order or as-built BOM. TODO: confirm this quote was
actually ordered and fitted before treating it as final.
\begin{itemize}
\item \textbf{Optional accessory:} HydroVolta's SonixED skid datasheet
  lists a \textbf{Cloud Multi-Parameter Water Analyzer (SUP-MPP1000)} as
  an available option --- a 7-inch touchscreen instrument measuring
  turbidity, pH, conductivity, dissolved oxygen and residual chlorine,
  with RS485/4--20\,mA output, for more detailed product-water quality
  monitoring than the standard sensor array.
\end{itemize}
TODO: confirm exact make/model per sensor type in the standard product.

\section{Operation}
Sensor readings feed the process-value displays described in
\cref{ch:operation} and the alarm limits in \cref{ch:appparams}.

\section{Maintenance}
\label{sec:sensorcal}
No formal calibration-to-reference procedure (e.g.\ pH 4/7/10 buffer
calibration) was found; what \emph{is} documented is a real bench
verification procedure used before commissioning any 2-wire 4--20\,mA
sensor:
\begin{procedure}
\item Connect the sensor per the electrical schematic up to the terminal
  strip --- \textbf{do not} connect to the PLC yet.
\item Insert a test resistor ($\sim$250\,$\Omega$, giving a safe 1--5\,V
  signal from the loop) and measure voltage with a voltage-only
  multimeter (e.g.\ Fluke~116).
\item For a pressure sensor (4\,mA~=~0\,bar): $\sim$0.88\,V across the test
  resistor confirms a normal 0\,bar signal before connecting to the PLC.
\item For a pH sensor (0--14\,pH, 4--20\,mA) in a 7\,pH reference solution:
  $\sim$2.64\,V across a 220\,$\Omega$ resistor (12\,mA at mid-scale) is
  expected.
\item If the sensor checks out on the bench but still misbehaves once
  installed, the fault is in the field wiring, not the sensor.
\end{procedure}
TODO: this is a wiring/signal-integrity check, not a full calibration
procedure --- a genuine calibration interval/method for pH, conductivity
and pressure sensors was not located.

\section{Troubleshooting}
\hvtablebegin{L L L}{4--20\,mA sensor fault table}{tab:sensorfault}
\hvhead{\textbf{Symptom (bench test)} & \textbf{Likely cause} & \textbf{Remedy}}\\\midrule
0\,V reading & No control voltage reaching sensor loop & Check DIN connector; confirm 24\,V\,DC between terminal and PE; check polarity; confirm supply is on\\
Reading below 4\,mA & Reversed polarity & Check polarity at the sensor connections\\
Reading stuck near 4\,mA & Sensor connection contamination & Check the connection is dry and clean\\
Reading above 20\,mA & Wiring fault / short & Check connections\\
No output signal at all (installed) & Cable break & Check loop continuity\\
Zero point has drifted & Overpressure limit exceeded, or wrong working temperature & Check against rated span and temperature range; recalibrate zero if within spec\\
Constant output despite pressure/process change & Mechanical fault in the sensor & Replace the sensor\\
Signal span varies over time & EMC interference & Shield the instrument/cable; remove the interference source\\
Signal span has dropped & Mechanical overload from an overpressure event & Replace the instrument\\
Fails on the bench too & Defective sensor & Replace the sensor\\
Passes bench test, fails installed & Field wiring fault & Inspect and repair wiring in the container\\
\bottomrule\end{xltabular}\par
The first four rows above are from the real HydroVolta bench verification
procedure; the following five are from the WIKA A-10 OEM manual's own
fault-finding section and apply directly to pressure sensors, and by the
same 2-wire 4--20\,mA principle, are a reasonable first check for the pH
and conductivity sensors as well.


\chapter{Control Cabinet, PLC and HMI}
\label{ch:cabinet}
\chaptersubtitle{The IP55 control cabinet, PLC I/O and operator touchscreen}
\chaptertoc

\section{Overview and function}
The control cabinet houses the PLC, I/O modules, VFDs and terminal wiring
that supervise the whole unit: a Weidm\"uller UC20-WL2000-AC controller
(CODESYS/u-OS based) coordinates distributed ADAM I/O modules
(ADAM-4024/4050/4017+/5000), the process VFDs (\cref{ch:pumps}), and the
dosing subsystem (\cref{ch:dosing}), with an optional touchscreen HMI for
local operation.

\hvfig{fig-control-cabinet.png}{Suggested figure: photograph of the open
  control cabinet with the PLC, I/O modules and terminal strips labeled.}
  {fig:control-cabinet}

\section{Key specifications}
\begin{itemize}
\item Cabinet ingress protection: IP55.
\item Controller: Weidm\"uller UC20-WL2000-AC (CODESYS/u-OS).
\item Distributed I/O: Advantech ADAM series. The \textbf{ADAM-3600-C2G}
  Wireless Intelligent RTU (also usable as the main controller) is fully
  datasheeted in HydroVolta's library: Cortex-A8 AM3352 @600\,MHz,
  256\,MB DDR3L RAM, RT-Linux~3.12; onboard 8AI/8DI/4DO plus a 4-slot I/O
  expansion bay (compatible modules: ADAM-3613/3617/3618/3624/3651/3656);
  1$\times$RS232/485 (DB9) + 2$\times$RS485 (terminal block), 2$\times$RJ-45
  10/100 Ethernet; protocols Modbus/TCP, DNP3~L2, IEC~104, MQTT; power
  10--30\,V\,DC (24\,V @ 5\,W typical); operating temperature
  $-$40 to 70\,\textdegree C; DIN-rail or wall mount. ADAM-4024 (analog
  output, used for dosing pump control), ADAM-4050/4017+ and the
  ADAM-5000-series backplane (analog I/O cards ADAM-5017P, ADAM-5024) are
  confirmed in real pilot I/O wiring documentation alongside the
  ADAM-3600. A more recent (2023) confirmed remote-I/O bill of materials
  uses \textbf{Weidm\"uller UR20} modular I/O instead: a
  UR20-FBC-MOD-TCP-V2 Modbus/TCP fieldbus coupler, one UR20-16DI-P
  (16-channel digital input), one UR20-16DO-P (16-channel digital
  output), and six UR20-4AI-UI-16 (4-channel universal analog input)
  modules on a DIN rail --- i.e.\ a Modbus/TCP-based I/O station with
  analog input as the dominant channel count, consistent with the many
  4--20\,mA sensors described in \cref{ch:sensors}. TODO: confirm whether
  the standard \salinbloc product uses the ADAM family, the Weidm\"uller
  UR20 family, or both depending on configuration --- both are confirmed
  purchases in HydroVolta's history, from different projects/years.
\item Control-box ordering options: \textbf{CB0} (PLC only), \textbf{CB1}
  (PLC~+ touchscreen HMI).
\item Cloud-connectivity tiers: \textbf{CL0} (local only), \textbf{CL1}
  (cloud monitoring), \textbf{CL2} (full cloud monitoring and control) ---
  see \cref{ch:telemetry}.
\item Incoming supply (pilot): 3$\times$400\,V\,AC (3P+N+PE), max.\ 25\,A;
  24\,V\,DC control bus.
\end{itemize}

\section{Operation}
The PLC drives the automatic sequencing described throughout Part~II
(pump speed, valve switching, dosing) and surfaces process values and
alarms to the HMI/dashboard (\cref{ch:operation}).

\section{Maintenance}
TODO: no cabinet-specific maintenance (dust/filter check on the IP55
enclosure, terminal torque check) was found.

\section{Troubleshooting}
\begin{itemize}
\item \textbf{ADAM-3600 timeout / cannot reconnect after uploading a new
  PLC program:} browse to the device's IP address, log in to its web UI
  (default password \texttt{00000000}), go to \emph{Configuration~$\to$~Reboot}.
  After reboot, retry the stop/connect sequence. Precondition: the device
  must respond to \texttt{ping} --- if it does not, the fault is network/
  power related, not a program-upload timeout.
\item \textbf{Modbus-RTU communication issue:} use qModMaster with a CH340
  USB-to-Modbus-RTU adapter to verify communication directly: confirm the
  assigned COM port, set RTU options to match the device (e.g.\
  9600\,baud, 8 data bits, 1 stop bit, no parity), set the correct slave
  address, use FC1/FC2/FC3, set Base Address to~0, then read a known
  register to confirm communication. The built-in bus monitor can inspect
  Tx/Rx traffic for deeper diagnosis.
\item \textbf{Analog signal fault (any 4--20\,mA device wired through this
  cabinet):} see the bench-test procedure in \cref{sec:sensorcal}.
\end{itemize}


\chapter{Tanks}
\label{ch:tanks}
\chaptersubtitle{Diluate, concentrate, electrolyte and drain tanks}
\chaptertoc

\section{Overview and function}
Overhead/skid-mounted tanks buffer the diluate, concentrate and
electrolyte streams (\cref{ch:overview}) and collect drainage from both
containers before it leaves the unit. In the pilot configuration these are
tagged TSA (a sensor/buffer tank), TB and TE (electrolyte) --- see
\cref{ch:sensors} for the level/conductivity instrumentation on each.

\hvfig{fig-tanks.png}{Suggested figure: photograph of the tank bank inside
  the treatment container, each tank labeled with its stream (diluate/
  concentrate/electrolyte/drain).}{fig:tanks}

\section{Key specifications}
No tank capacity/material specification exists yet for a finalized
standard \salinbloc bill of materials, but two real, recent tank orders
placed by HydroVolta give a concrete sense of the sizes used:
\begin{itemize}
\item Rotterdam Plastics order (Oct.\ 2025): one 800-litre vertical
  polyethylene tank (closed outlet, unvented lid) and one 450-litre
  polyethylene tank (open outlet, vented lid).
\item Boralit/Galo Tanks order: two 1000-litre rectangular polyethylene
  tanks (RT1000RH~ZW series), supplied with PVC coupling sets.
\end{itemize}
TODO: which physical tank (diluate/feed, concentrate, electrolyte, or
drain) each of the above corresponds to was not confirmed in the order
documents --- do not assume a mapping without checking the as-built P\&ID
for the specific unit. The pilot configuration referenced elsewhere in
this manual used a 1000\,l IBC as an auxiliary/drain tank.

\section{Operation}
Tank levels are monitored continuously (\plctag{TSA\_Level},
\plctag{TB\_Level}, \plctag{TE\_Level}) with high/low alarm limits ---
see \cref{ch:appparams}.

\section{Maintenance}
TODO: no tank cleaning/inspection schedule was found.

\section{Troubleshooting}
High/low level alarms are confirmed for every tank:
\plctag{HL\_Alarm\_TD}/\plctag{LL\_Alarm\_TD} (diluate),
\plctag{HL\_Alarm\_TC}/\plctag{LL\_Alarm\_TC} (concentrate),
\plctag{HL\_Alarm\_TB}/\plctag{LL\_Alarm\_TB},
\plctag{HL\_Alarm\_TSA}/\plctag{LL\_Alarm\_TSA},
\plctag{HL\_Alarm\_TE}/\plctag{LL\_Alarm\_TE}, plus a level alarm
\plctag{M050\_Level\_Alarm}. On any level alarm, check the relevant level
sensor per the bench-test procedure in \cref{sec:sensorcal} before
assuming a process fault.


\chapter{Telemetry and Remote Connectivity}
\label{ch:telemetry}
\chaptersubtitle{Router, remote-access gateway and cloud monitoring}
\chaptertoc

\section{Overview and function}
A cellular router and cloud dashboard give operators and HydroVolta
service remote visibility into the unit's process values and alarms,
independent of the local HMI (\cref{ch:cabinet}).

\hvfig{fig-telemetry-router.png}{Suggested figure: photograph of the
  router/gateway mounted in the cabinet, with antenna and SIM slot
  visible.}{fig:telemetry-router}

\section{Key specifications}
\begin{itemize}
\item \textbf{Confirmed deployed remote access:} real OpenVPN client
  configuration files exist in HydroVolta's engineering library for at
  least two pilot containers (Farys, BASF/CAGECE demo), following a
  router-based \enquote{meeting point} VPN architecture (Weidm\"uller
  IE-SR-2GT-LAN / IE-SR-2GT-UMTS routers, per the vendor tech note on
  file) --- this is the one telemetry element confirmed as \emph{actually
  used}, not just evaluated.
\item \textbf{Confirmed deployed dashboard:} \textbf{ThingsBoard}, with a
  real MQTT parameter list and full tag list captured from the BASF pilot
  deployment.
\item A real 2023 requirements-workshop record (source:
  \texttt{08\_IT/00\_OQuila progress}) describes an \textbf{actually
  running 4-container pilot fleet} at that time: \textbf{Advantech
  ADAM-3600} connecting directly to \textbf{Azure IoT~Hub}, piped through
  \textbf{Azure Stream Analytics} into a SQL database and visualized in
  \textbf{Power~BI}; connectivity via a \textbf{4G Teltonika router}, with
  HydroVolta noting a possible future switch to Ewon \enquote{for a
  better SLA} due to past Teltonika stability issues. Each container sent
  a telemetry message roughly every 10\,seconds ($\approx$8{,}000
  messages/container/day: connectivity, pressure, flow, temperature,
  energy). Parameter changes at that time were only possible via VPN, not
  from the cloud dashboard.
\item \textbf{OQuila} was a third-party IoT consultancy HydroVolta held a
  2023 requirements workshop with, to evaluate a longer-term telemetry
  platform (long-term storage, a graphical P\&ID front-end, configurable
  alarm thresholds, role-based access). HydroVolta's own workshop notes
  state ThingsBoard was assessed as \enquote{ok} for the dashboard need on
  its own. \textbf{No document confirms OQuila's proposal was ultimately
  adopted or built} --- treat the Azure~IoT~Hub/Power~BI pipeline above
  and the ThingsBoard/ Weidm\"uller-OpenVPN pipeline elsewhere in this
  manual as two telemetry architectures documented at different points in
  HydroVolta's history, not necessarily the same one. TODO: confirm which
  architecture (if either, unchanged) is used on the current standard
  \salinbloc product.
\item Controller reached over the network: Weidm\"uller UC20-WL2000-AC
  (\cref{ch:cabinet}).
\end{itemize}

\section{Operation}
Process values available on the cloud dashboard mirror those on the local
HMI (\cref{ch:operation}). TODO: no document describes ThingsBoard
dashboard login/navigation for an end customer.

\section{Maintenance}
TODO: no router/SIM maintenance procedure (data plan renewal, firmware
update schedule) was found.

\section{Troubleshooting}
\begin{itemize}
\item \textbf{Router unreachable:} confirm the device is powered and check
  its cellular signal/SIM status before assuming a software fault.
\item \textbf{Cloud dashboard shows stale data:} check the router's
  connection to the controller first; if the controller itself is
  reachable locally but not remotely, the fault is likely in the
  router/network path rather than the controller.
\end{itemize}
See also the ADAM-3600 and Modbus-RTU troubleshooting steps in
\cref{ch:cabinet}, which apply to field-device connectivity issues more
broadly.


%% ---- Part III: Closing general chapters -----------------------------------

\chapter{Shutdown, storage and disposal}
\label{ch:shutdown}
\chaptertoc

\section{Short shutdown}
Use \textbf{Controlled Shutdown} (Mode~4) from the Mode Control HMI
screen (\cref{ch:operation}) rather than simply removing power, so pumps
and valves stop in a defined sequence. TODO: no document describes a
written de-energization procedure beyond mode selection; the
incoming-supply breaker layout (main isolator, MCBs, RCDs) documented on
the pretreatment electrical schematic would be the basis for one if the
unit is to be fully powered down (as opposed to left in Controlled
Shutdown/standby).

\section{Long-term storage}
No storage/preservation procedure specific to the whole \salinbloc unit
was found, but the UF membrane manufacturer's preservation procedure
(\cref{ch:pretreatment}) is directly applicable to the pretreatment stage
and is the best-documented reference available: soak the UF module in a
glycerol/sodium-metabisulfite preservation solution before an extended
shutdown, and flush with $\geq$14\,m\textsuperscript{3} of feed water for
at least 4\,hours before restart. TODO: no equivalent preservation
guidance (membrane wetting, antifouling flush, dosing-system
winterization) was found for the \sonixed stack itself or the dosing
system --- for the dosing system, at minimum follow the chemical storage
guidance in \cref{ch:dosing} (15--25\,\textdegree C, tightly closed
original containers) for the duration of the shutdown.

\section{Restarting after storage}
TODO: no start-up/recommissioning-after-storage procedure was found; treat
as a fresh commissioning (\cref{ch:commissioning}) until a dedicated
procedure exists.

\section{Decommissioning}
TODO: no decommissioning procedure was found.

\section{Disposal and recycling}
TODO: no disposal/recycling procedure or waste-stream document was found.
Component datasheets (e.g.\ dosing pumps, \cref{ch:dosing}) may carry
manufacturer disposal notes that have not yet been reviewed for this
section.


\appendix

\chapter{Parameter tables}
\label{ch:appparams}
\chaptertoc

\section{Default setpoints}
\hvtablebegin{L R{22mm}}{System performance defaults (SonixED skid v2025-A7)}{tab:defaultsetpoints}
\hvhead{\textbf{Parameter} & \textbf{Value}}\\\midrule
Flowrate per stack & 0.5--10\,m\textsuperscript{3}/h\\
System flow per skid & up to 30\,m\textsuperscript{3}/h\\
Recovery rate & up to 98\%\\
Operating pressure & 0.5--2.5\,bar\\
Temperature range & 5--55\,\textdegree C\\
Specific energy & 0.6--2.0\,kWh/m\textsuperscript{3}\\
Skid dimensions (l$\times$b$\times$h) & 3040$\times$1340$\times$2020\,mm\\
Incoming line voltage & 3$\times$400\,V\,AC\\
Max.\ current & 25\,A\\
Valve voltage & 24\,V\,DC\\
Control voltage (DC bus) & 24\,V\,DC\\
Dosing SET~1 (min) & 4.0\,mA / 0 impulses/min\\
Dosing SET~2 (max) & 20.0\,mA / 180 impulses/min\\
\bottomrule\end{xltabular}\par

\hvtablebegin{L R{20mm} R{20mm}}{\salinbloc container throughput by format (source: HydroVolta brochure)}{tab:containerthroughput}
\hvhead{\textbf{Parameter} & \textbf{20\,ft unit} & \textbf{40\,ft unit}}\\\midrule
Flow rate & 2--30\,m\textsuperscript{3}/h & 15--60\,m\textsuperscript{3}/h\\
Approx.\ daily throughput (22\,h/day operation) & 44--660\,m\textsuperscript{3}/day & 330--1{,}320\,m\textsuperscript{3}/day\\
\bottomrule\end{xltabular}\par
Feed range for both formats: brackish groundwater, 1{,}500--20{,}000\,ppm
TDS. TODO: no container weight figure was found in any source reviewed,
including the container general-arrangement drawings themselves (their
title blocks carry no weight/dimension text field, and any dimension
lines are vector geometry, not extractable text --- a human must open
these drawings visually). Obtain the weight from the container
manufacturer or a loaded-unit weighbridge reading before publishing
lifting/foundation guidance.

\section{Site-specific settings}
TODO: no site-specific commissioning/setpoint record was found; use this
table to capture the as-commissioned values for a specific installation.

\hvtablebegin{L L}{Site-specific settings template}{tab:sitesettings}
\hvhead{\textbf{Parameter} & \textbf{Value}}\\\midrule
 & \\
 & \\
\bottomrule\end{xltabular}\par

\section{Alarm limits}
The P\&ID (\cref{ch:appdrawings}) shows instrument spans used as default
alarm spans at each skid position: pressure 0--2.5\,bar, conductivity
0--200\,mS/cm, with high/low alarm tags (PSA, CEA) at every position.
HydroVolta's ThingsBoard warning-code table (a real, deployed alarm
definition list, source: \texttt{08\_IT/01\_Thingsboard/Warning table})
gives concrete numeric examples for several alarm classes:

\hvtablebegin{L L L}{Representative alarm/warning codes (source: HydroVolta ThingsBoard warning table)}{tab:alarmcodes}
\hvhead{\textbf{Code} & \textbf{Class} & \textbf{Condition}}\\\midrule
W001 & General & HV-Skid not powered\\
W002 & General & Ventilation off\\
W011--W013 & Chemical & Antiscalant / H\textsubscript{2}SO\textsubscript{4} / HCl tank empty\\
W101--W110 & Level & Main diluate/concentrate/electrolyte/breaking/client tank level low or high (e.g.\ diluate below 40\,cm or above 90\,cm; electrolyte below 20\,cm or above 80\,cm)\\
W111--W118 & Level & Buffer tank low-low / high level\\
W201--W214 & Electrical & Stack stage 1--7 voltage below or above system specification\\
W300 & Flow & Diluate flow below 500\,l/hr\\
W301 & Flow & Concentrate flow below 500\,l/hr\\
\bottomrule\end{xltabular}\par
TODO: this table only samples the documented code ranges; the full
warning table has $\sim$57 defined codes and should be transcribed in
full by HydroVolta engineering, together with the exact level/pressure
setpoints per tank size, once the standard \salinbloc tank configuration
(\cref{ch:tanks}) is finalized --- these codes come from a specific pilot
project's tank/stage configuration and the numeric thresholds may not
transfer directly. See the alarm tag list in
\cref{ch:operation,ch:tanks,ch:dosing} for additional alarms.


\chapter{Drawings}
\label{ch:appdrawings}
\chaptertoc

\section{P\&ID}
Piping \& instrumentation diagram for the containerized BASF pilot water
treatment system: drawing no.\ SWHC-001/00, rev.\ 0.98, dated 11-11-2022,
\enquote{BASF WATER Container} (Grant Agreement 969116~--- SonixED,
H2020-EIC-SMEInst). Shows pretreatment filtration skids (SP-01--SP-4),
dosing units for antiscalant/NaOCl/NaOH/HCl/H\textsubscript{2}SO\textsubscript{4}
(CIP), tanks, and instrumentation with alarm tags across the treatment
train.

\hvfig{fig-pid.png}{Insert the full P\&ID drawing here (source: SWHC-001/00,
  rev.\,0.98) once confirmed representative of the standard \salinbloc
  product, or the product-specific P\&ID once available.}{fig:pid}

HydroVolta's standardization library includes an ANSI/ISA-5.1 P\&ID
symbol legend, which is the reference standard the symbols on this and
other project-specific P\&IDs (BASF, Saudi, Australia demo, Morocco) are
drawn against.

\section{Electrical schematic}
Circuit diagram set \enquote{Pretreatment EP} for the pretreatment
container, document series C1/I1/I2/E1--E17 (file \texttt{31-01-2023
Pretreatment.dsn}). Covers the 3-phase incoming supply and breaker/RCD
protection (sheet~E1), 24\,V\,DC power/control distribution (E2--E3), PLC
I/O wiring (E4--E5), terminal wiring (E6--E8), valve/sensor junction boxes
(E9--E13), VFD control (E14), sensor/valve wiring standards (E15), dosing
pump wiring (E16), and Modbus COM2 wiring (E17).

\hvfig{fig-electrical-schematic.png}{Insert the relevant electrical
  schematic sheet(s) here (source: \texttt{31-01-2023
  Pretreatment.dsn}, sheets E1--E17) once confirmed representative of the
  standard \salinbloc product.}{fig:electrical-schematic}

\section{Container layout}
20\,ft general-arrangement drawing \enquote{2023\_08\_21 SonixED container
BASF\_2308CA1} (6058\,$\times$\,2438\,$\times$\,2896\,mm), and companion
skid top-view layout \enquote{2023\_08\_22 BASF Skids Top
View\_2308CA2} --- see \cref{ch:overview} for the summary description.

\hvfig{fig-container-layout-drawing.png}{Insert drawings 2308CA1 (general
  arrangement) and 2308CA2 (skid top view) here, or the standard-product
  equivalents once available.}{fig:container-layout-drawing}


\chapter{Certificates and datasheets}
\label{ch:appcert}
\chaptertoc

\section{Declaration of conformity}
TODO: only a blank EU Machinery Directive DoC template was found
(\enquote{EU declaration of conformity-en}, Ref.\ Ares(2015)1600946), not a
signed, product-specific Declaration of Conformity --- HydroVolta's
\texttt{02\_DoC's} CE-marking folder is currently empty. A signed DoC must
be inserted here before this manual is issued to a customer.

HydroVolta's legal/legislation library holds reference texts for the EU
directives and regulations most likely to apply to a CE-marked
\salinbloc unit, which a future DoC should be checked against: the
Machinery Directive (2006/42/EC), EMC Directive (2014/30/EU), Low Voltage
Directive (2014/35/EU), RoHS, WEEE, the EU Drinking Water Directive
(2020/2184) if a potable-water claim is made, and REACH/CLP for the
dosing chemicals (\cref{ch:dosing}). A component-level DoC exists for at
least one purchased part (the WIKA A-10 pressure sensor, \cref{ch:sensors})
--- this is normal for a CE-marked assembly (component conformity feeds
the assembly-level DoC) but is not itself a substitute for the unit's own
Declaration of Conformity.

\section{Component datasheets}
The \sonixed Skid System datasheet (Model SONIXED-SKID, Version 2025-A7)
is the primary component reference used throughout Part~II of this manual.
Individual component-level datasheets (I/O modules, dosing pumps,
pumps, pneumatics) exist in HydroVolta's engineering library and should be
attached here per shipped configuration.

\section{Regulatory approvals (drinking water)}
No UK DWI Regulation~31 approval and no NSF/ANSI listing (e.g.\
NSF/ANSI~61 or~372) exists yet for \salinbloc/\sonixed --- HydroVolta's
datasheet library holds only generic reference/guidance material on how
those approval processes work, not an application or approval record (see
\cref{sec:intendeduse} in \cref{ch:safety}). TODO: pursue formal approval
before offering the unit for potable-water use in a regulated market, or
clearly disclaim potable-water suitability until then.

\section{Water analysis report}
TODO: no formal water-analysis lab report/certificate of analysis was
found for the standard product; the closest material found (third-party
pretreatment-component water-quality questionnaires from earlier pilot
projects) is not a substitute. A per-unit water analysis report should be
attached here once available --- see also the "Find your unit's quality
report" lookup on the online container guide page (\cref{ch:about}).


\chapter{Glossary}
\label{ch:appglossary}

\abbreviationsbegin
\abbr{TDS}{Total Dissolved Solids --- the concentration of dissolved minerals/salts in water.}
\abbr{\sonixed}{HydroVolta's electrodialysis-based desalination technology/product line.}
\abbr{ED}{Electrodialysis --- a membrane process using an electric field to move ions through selective membranes.}
\abbr{EDR}{Electrodialysis Reversal --- periodically reversing the applied DC field's polarity to clean the membranes in place and stabilize performance.}
\abbr{TMP}{Transmembrane Pressure --- the pressure difference across a filtration membrane; a rising TMP indicates fouling and triggers cleaning.}
\abbr{CEB}{Chemically Enhanced Backwash --- a backwash cycle that includes a chemical cleaning soak, used on the ultrafiltration stage.}
\abbr{VCA}{Veiligheid, Gezondheid en Milieu Checklist Aannemers --- a Benelux contractor health/safety/environment competency certification.}
\abbr{Stack}{The assembly of alternating ion-exchange membranes, electrodes and spacers that performs electrodialysis.}
\abbr{Skid}{A pre-assembled equipment module (pumps, valves, instrumentation) on a common frame.}
\abbr{Diluate}{The desalinated water stream leaving the stack.}
\abbr{Concentrate}{The brine/reject stream leaving the stack, carrying the removed ions.}
\abbr{Electrolyte}{A separate circulating stream feeding the stack's electrode compartments.}
\abbr{Ion-exchange membrane}{A selective membrane that passes either cations or anions, blocking the other.}
\abbr{PLC}{Programmable Logic Controller --- the industrial controller running the unit's automatic logic.}
\abbr{HMI}{Human-Machine Interface --- the operator touchscreen/display.}
\abbr{VFD}{Variable-Frequency Drive --- controls pump/motor speed.}
\abbr{RTU}{Remote Terminal Unit --- a networked I/O/communication device (e.g.\ the ADAM-3600).}
\abbr{CIP}{Clean-In-Place --- an automated cleaning cycle using dosed chemicals.}
\abbr{P\&ID}{Piping \& Instrumentation Diagram.}
\abbr{DoC}{Declaration of Conformity --- the manufacturer's statement of regulatory compliance (e.g.\ CE marking).}
\abbr{PPE}{Personal Protective Equipment.}
\abbr{E-stop}{Emergency stop --- a manually operated device that immediately halts machine operation.}
\bottomrule\end{xltabular}\par

\hvbackcover

\end{document}
